% !TEX program = xelatex
% 任务规划算法体系 - 教材
% 编译方式: xelatex main.tex (运行两次以生成目录)

\documentclass[a4paper,12pt,openany]{ctexbook}

% ============ 基础宏包 ============
\usepackage[top=2.5cm,bottom=2.5cm,left=2.5cm,right=2.5cm]{geometry}
\usepackage{amsmath,amssymb,amsthm}
\usepackage{graphicx}
\usepackage{booktabs}
\usepackage{longtable}
\usepackage{tabularx}
\usepackage{multirow}
\usepackage{enumitem}
\usepackage{tikz}
\usetikzlibrary{shapes,arrows,positioning,calc,fit,backgrounds,decorations.pathreplacing}
\usepackage{algorithm}
\usepackage{algpseudocode}
\usepackage{listings}
\usepackage{xcolor}
\usepackage{tcolorbox}
\tcbuselibrary{skins,breakable,theorems}
\usepackage{fancyhdr}
\usepackage{titlesec}

% hyperref 放在最后加载
\usepackage[colorlinks=true,linkcolor=blue!70!black,citecolor=green!50!black,urlcolor=blue!60!black,bookmarks=true,bookmarksnumbered=true]{hyperref}

% ============ 颜色定义 ============
\definecolor{ugcolor}{RGB}{0,102,204}       % 本科生 - 蓝色
\definecolor{gradcolor}{RGB}{204,51,0}      % 研究生 - 红色
\definecolor{selfcolor}{RGB}{0,153,76}      % 自学拓展 - 绿色
\definecolor{codebg}{RGB}{245,245,245}
\definecolor{codeframe}{RGB}{200,200,200}
\definecolor{partcolor}{RGB}{0,51,102}

% ============ 受众标记命令 ============
% 本科生适用
\newcommand{\ug}{\texorpdfstring{\textcolor{ugcolor}{\textbf{[本]}}}{[本]}}
% 研究生扩展
\newcommand{\grad}{\texorpdfstring{\textcolor{gradcolor}{\textbf{[研]}}}{[研]}}
% 自学拓展
\newcommand{\self}{\texorpdfstring{\textcolor{selfcolor}{\textbf{[拓]}}}{[拓]}}

% ============ 定理环境 ============
\theoremstyle{definition}
\newtheorem{definition}{定义}[chapter]
\newtheorem{example}{例}[chapter]
\newtheorem{remark}{注}[chapter]
\theoremstyle{plain}
\newtheorem{theorem}{定理}[chapter]
\newtheorem{proposition}{命题}[chapter]
\newtheorem{lemma}{引理}[chapter]

% ============ tcolorbox 环境 ============
\newtcolorbox{keypoint}[1][]{
  colback=blue!5!white,
  colframe=blue!60!black,
  title={\textbf{要点}},
  fonttitle=\bfseries,
  breakable,
  #1
}

\newtcolorbox{algbox}[1][]{
  colback=gray!5!white,
  colframe=gray!60!black,
  title={\textbf{算法概要}},
  fonttitle=\bfseries,
  breakable,
  #1
}

\newtcolorbox{appbox}[1][]{
  colback=green!5!white,
  colframe=green!60!black,
  title={\textbf{应用案例}},
  fonttitle=\bfseries,
  breakable,
  #1
}

\newtcolorbox{levelbox}[1][]{
  colback=orange!5!white,
  colframe=orange!60!black,
  title={\textbf{学习建议}},
  fonttitle=\bfseries,
  breakable,
  #1
}

% ============ 代码环境 ============
\lstset{
  basicstyle=\ttfamily\small,
  backgroundcolor=\color{codebg},
  frame=single,
  rulecolor=\color{codeframe},
  breaklines=true,
  numbers=left,
  numberstyle=\tiny\color{gray},
  keywordstyle=\color{blue!70!black}\bfseries,
  commentstyle=\color{green!50!black},
  stringstyle=\color{red!60!black},
  showstringspaces=false
}

% ============ 页眉页脚 ============
\pagestyle{fancy}
\fancyhf{}
\fancyhead[LE]{\leftmark}
\fancyhead[RO]{\rightmark}
\fancyfoot[C]{\thepage}
\renewcommand{\headrulewidth}{0.4pt}

% ============ 章节标题格式 ============
\ctexset{
  part/format = \Huge\bfseries\centering\color{partcolor},
  chapter/format = \LARGE\bfseries\color{blue!70!black},
  section/format = \Large\bfseries,
  subsection/format = \large\bfseries
}

% ============ 算法环境中文化 ============
\floatname{algorithm}{算法}
\renewcommand{\algorithmicrequire}{\textbf{输入:}}
\renewcommand{\algorithmicensure}{\textbf{输出:}}

% ============ 文档信息 ============
\title{{\Huge\bfseries 《任务规划》之算法部分}\\[1em]
{\Large 从经典方法到智能规划}\\[2em]
{\normalsize\color{gray} 教学参考教材}}
\author{}
\date{\today}

\begin{document}

% ============ 封面 ============
\maketitle

% ============ 受众标记说明 ============
\thispagestyle{empty}
\vspace*{3cm}
\begin{center}
{\Large\bfseries 本书受众标记说明}
\end{center}
\vspace{1cm}

\noindent 本书内容覆盖面较广，为方便不同层次的读者使用，采用以下标记区分内容的适用层次：

\vspace{1em}

\begin{center}
\begin{tabular}{clp{8cm}}
\toprule
\textbf{标记} & \textbf{适用层次} & \textbf{说明} \\
\midrule
\textcolor{ugcolor}{\textbf{[本]}} & 本科生教学 & 核心基础内容，建议本科生课堂讲授和学习。掌握这些内容即可建立完整的任务规划知识体系。 \\[0.5em]
\textcolor{gradcolor}{\textbf{[研]}} & 研究生扩展 & 进阶内容，适合研究生课程或本科高年级选修。涉及前沿算法和深入理论分析。 \\[0.5em]
\textcolor{selfcolor}{\textbf{[拓]}} & 自学拓展 & 补充阅读材料，适合感兴趣的读者自学。包含最新研究进展和特殊应用领域。 \\
\bottomrule
\end{tabular}
\end{center}

\vspace{2em}
\noindent\textbf{建议学习路线：}
\begin{itemize}[leftmargin=2em]
  \item \textbf{本科生（32--48学时）：}重点学习标记为 \ug{} 的章节，建立扎实的基础。
  \item \textbf{研究生（48--64学时）：}在 \ug{} 内容基础上，深入学习 \grad{} 标记的内容。
  \item \textbf{自学参考：}可根据兴趣选择 \self{} 标记的内容进行拓展阅读。
\end{itemize}

\clearpage

% ============ 目录 ============
\tableofcontents

% ============ 前言 ============
\chapter*{前言}
\addcontentsline{toc}{chapter}{前言}

任务规划（Task Planning）是人工智能领域的核心研究方向之一，其目标是让智能体能够自主地制定行动方案以达成预设目标。从早期的GPS（General Problem Solver）到当今基于大语言模型的智能规划系统，任务规划技术经历了数十年的发展，形成了丰富而完整的算法体系。

本书旨在系统梳理任务规划领域的主要算法，建立清晰的方法论体系。全书按照算法的技术路线和演化脉络，分为六个篇章：

\begin{itemize}[leftmargin=2em]
  \item \textbf{第一篇：经典规划方法}——介绍基于符号推理和搜索的传统规划算法，这是整个领域的理论基石。
  \item \textbf{第二篇：元启发式与优化方法}——介绍基于群体智能和进化计算的规划求解方法。
  \item \textbf{第三篇：强化学习与规划}——探讨基于试错学习的规划方法，包括基于模型和无模型的方法。
  \item \textbf{第四篇：深度学习与规划}——介绍利用深度神经网络提升规划能力的现代方法。
  \item \textbf{第五篇：大语言模型与规划}——聚焦LLM驱动的最新规划技术，包括具身智能和语言智能体。
  \item \textbf{第六篇：多智能体规划}——讨论多个智能体协同规划的理论与方法。
\end{itemize}

每个篇章中，我们既关注算法的核心思想和理论分析，也注重其在交通调度、通信编码、军事指挥等实际领域的应用。书中采用受众标记系统，帮助不同层次的读者高效地选择学习内容。

近年来，国内学者在任务规划相关领域取得了丰富的研究成果。在元启发式优化方面，段海滨、焦李成、陈杰等学者的工作具有重要影响；在强化学习方面，高阳、赵冬斌、刘全等的综述论文为中文读者提供了系统的学习资料；在无人机集群任务规划方面，邢立宁、贾高伟等的综述系统梳理了该领域的进展；在大语言模型与规划方面，秦龙等提出了面向复杂任务的自主规划框架。本书在介绍国际前沿的同时，也充分融入了国内学者的研究贡献和国内应用场景，力求为中文读者提供兼具国际视野与本土特色的参考教材。

\vspace{2em}
\hfill 编者

\hfill \today

% =============================================
% 第一篇：经典规划方法
% =============================================
\part{经典规划方法}

% 第1章：任务规划基础与状态空间搜索
\chapter{任务规划基础与状态空间搜索 \ug}

\begin{levelbox}
本章为全书基础，所有层次读者均需掌握。内容涵盖任务规划的基本概念、形式化描述和状态空间搜索方法，是理解后续所有算法的前提。
\end{levelbox}

% =============================================
\section{任务规划概述}

\subsection{什么是任务规划}

任务规划（Task Planning），又称自动规划（Automated Planning），是人工智能的重要分支。其核心问题可以描述为：给定一个初始状态、一组可执行的动作和一个目标条件，自动生成一个动作序列（即\textbf{规划}），使得从初始状态出发依次执行这些动作后能够到达满足目标条件的状态。

\begin{remark}
关于任务规划/自动规划的系统性学习资源，国内外均有优秀的教材。国际经典教材首推Ghallab等著的\emph{Automated Planning: Theory and Practice}（国内有殷建平等翻译的中文版《自动规划：理论和实践》）。国内人工智能教材如蔡自兴《人工智能及其应用》（第7版）和吴飞《人工智能引论》（``101计划''核心教材）均有专门章节讨论智能规划。邢立宁和陈英武在《火力与指挥控制》上发表的``任务规划系统研究及发展''综述了任务规划系统的整体发展脉络。
\end{remark}

\begin{definition}[经典规划问题]
一个经典规划问题可以形式化为四元组 $\Pi = \langle \mathcal{F}, \mathcal{A}, I, G \rangle$，其中：
\begin{itemize}
  \item $\mathcal{F}$：有限的命题变量集合（状态变量）；
  \item $\mathcal{A}$：有限的动作集合，每个动作 $a \in \mathcal{A}$ 由前置条件 $\mathrm{pre}(a)$、添加效果 $\mathrm{add}(a)$ 和删除效果 $\mathrm{del}(a)$ 描述；
  \item $I \subseteq \mathcal{F}$：初始状态；
  \item $G \subseteq \mathcal{F}$：目标条件。
\end{itemize}
\end{definition}

\subsection{规划与搜索的关系}

任务规划本质上是一种搜索问题。状态空间构成了一个有向图，其中节点是状态，边是动作。规划的过程就是在这个图中寻找从初始状态到目标状态的路径。

\subsection{PDDL：规划领域定义语言}

PDDL（Planning Domain Definition Language）是规划领域的标准建模语言，由1998年的第一届国际规划竞赛（IPC）引入。PDDL将规划问题分为\textbf{领域描述}（Domain）和\textbf{问题描述}（Problem）两部分。

\begin{example}[积木世界的PDDL描述]
积木世界（Blocksworld）是经典的规划基准问题。假设有三个积木A、B、C，初始状态为A在桌上、B在A上、C在B上（即C-B-A-桌），目标是将它们排列为A-B-C-桌。

领域文件定义了四个基本动作：拿起（pick-up）、放下（put-down）、堆叠（stack）、卸载（unstack）。问题文件则指定具体的初始状态和目标状态。
\end{example}

% =============================================
\section{无信息搜索方法 \ug}

无信息搜索（Uninformed Search）又称盲目搜索，不使用问题的特定知识来引导搜索方向。

\subsection{宽度优先搜索（BFS）}

宽度优先搜索按层次逐层扩展节点，保证找到最短路径（步数最少的规划）。

\begin{algbox}[title={\textbf{BFS基本思想}}]
\begin{enumerate}
  \item 将初始状态放入队列；
  \item 取出队首状态，检查是否满足目标；
  \item 若不满足，生成所有后继状态并加入队尾；
  \item 重复直到找到目标或队列为空。
\end{enumerate}
\textbf{完备性：}是\quad \textbf{最优性：}是（步数最优）\quad \textbf{时间复杂度：}$O(b^d)$\quad \textbf{空间复杂度：}$O(b^d)$
\end{algbox}

\subsection{深度优先搜索（DFS）}

深度优先搜索沿一条路径深入到底再回溯，空间效率高但不保证最优性。

\subsection{迭代加深搜索（IDS）}

迭代加深搜索结合了BFS的完备性和DFS的空间效率，是无信息搜索中最实用的方法之一。它以递增的深度限制反复执行深度有限搜索。

\subsection{双向搜索}

双向搜索同时从初始状态和目标状态出发进行搜索，当两个搜索前沿相遇时找到解。理论上可将搜索复杂度从$O(b^d)$降低到$O(b^{d/2})$。

% =============================================
\section{启发式搜索方法 \ug}

启发式搜索（Heuristic Search）利用问题特定的估价函数引导搜索方向，是现代规划器的核心技术。

\subsection{贪婪最佳优先搜索（GBFS）}

贪婪最佳优先搜索总是扩展启发值$h(n)$最小的节点。速度快但不保证最优性。

\subsection{A*搜索}

\begin{definition}[A*搜索]
A*搜索使用评价函数 $f(n) = g(n) + h(n)$，其中$g(n)$是从初始状态到$n$的实际代价，$h(n)$是从$n$到目标的启发式估计。当$h(n)$是可接纳的（即不超过实际代价）时，A*保证找到最优解。
\end{definition}

\begin{keypoint}
A*搜索的关键性质：
\begin{itemize}
  \item \textbf{可接纳性（Admissibility）：}若$h(n) \leq h^*(n)$对所有$n$成立（$h^*$为实际代价），则A*找到最优解。
  \item \textbf{一致性（Consistency）：}若$h(n) \leq c(n,a,n') + h(n')$，则A*不需要重新打开已关闭的节点。
  \item 在相同启发函数下，A*扩展的节点数不超过任何其他最优搜索算法。
\end{itemize}
\end{keypoint}

\subsection{IDA*搜索}

IDA*（Iterative Deepening A*）将迭代加深与A*结合，使用$f$值作为深度限制的阈值。空间复杂度仅为$O(bd)$，适合内存受限的场景。

\subsection{加权A*搜索 \grad}

加权A*使用$f(n) = g(n) + w \cdot h(n)$（$w > 1$），牺牲最优性换取搜索速度。找到的解的代价不超过最优解的$w$倍。

% =============================================
\section{规划中的搜索策略}

\subsection{前向搜索（前进状态空间搜索）}

从初始状态出发，通过应用动作逐步向目标推进。大多数现代规划器采用这种方式。

\subsection{后向搜索（回归搜索）}

从目标条件出发，反向寻找能够达成目标的动作序列。在某些问题中效率更高。

\subsection{搜索中的剪枝技术}

\begin{itemize}
  \item \textbf{重复状态检测：}使用哈希表记录已访问状态，避免重复探索。
  \item \textbf{有用动作剪枝（Helpful Actions）：}FF规划器引入的技术，只考虑"有用"的动作以减少分支因子。
  \item \textbf{对称性消除：}识别和排除对称等价的搜索分支。
\end{itemize}

% =============================================
\section{本章小结与习题}

\subsection*{本章小结}
\begin{itemize}
  \item 任务规划的核心是在状态空间中搜索从初始状态到目标状态的动作序列。
  \item PDDL是规划问题的标准描述语言。
  \item 无信息搜索提供基本保障，启发式搜索是实用规划器的基础。
  \item A*搜索在可接纳启发函数下保证最优性。
\end{itemize}

\begin{appbox}[title={\textbf{应用：中国高铁调度任务规划}}]
中国拥有世界最大规模的高速铁路网络（超过4.5万公里），列车运行图的编排是一个典型的大规模任务规划问题。需要在满足安全间隔、车站容量、动车组周转等约束下，为数千列列车规划运行路径和时刻。该问题可以建模为带时间窗和资源约束的经典规划问题，其规模和复杂度远超一般的学术基准问题。
\end{appbox}

\subsection*{思考与练习}
\begin{enumerate}
  \item 用PDDL描述一个简单的物流配送问题。
  \item 比较BFS和DFS在积木世界问题上的性能差异。
  \item 证明：若启发函数$h$是一致的，则它也是可接纳的。
  \item 在一个$4 \times 4$的网格世界中，比较A*和贪婪最佳优先搜索的节点扩展数量。
  \item 分析中国高铁调度问题中，搜索空间的状态数量大约是什么量级？为什么需要启发式搜索？
\end{enumerate}

% 第2章：启发式函数设计
\chapter{启发式函数设计 \ug}

\begin{levelbox}
启发式函数是现代规划器的核心组件。本章1--3节为本科生核心内容，第4--5节适合研究生深入学习。
\end{levelbox}

% =============================================
\section{松弛启发式 \ug}

\subsection{删除松弛（Delete Relaxation）}

删除松弛是规划中最重要的问题简化技术。其核心思想是：忽略所有动作的删除效果，使得一旦某个命题变为真就永远保持为真。

\begin{definition}[删除松弛问题]
给定规划问题 $\Pi = \langle \mathcal{F}, \mathcal{A}, I, G \rangle$，其删除松弛问题 $\Pi^+ = \langle \mathcal{F}, \mathcal{A}^+, I, G \rangle$，其中每个动作 $a^+ \in \mathcal{A}^+$ 的删除效果为空：$\mathrm{del}(a^+) = \emptyset$。
\end{definition}

\subsection{$h^{\max}$启发式}

$h^{\max}$估计达成目标中最难实现的单个子目标的代价：
$$h^{\max}(s) = \max_{g \in G} h^{\max}_1(s, g)$$
其中$h^{\max}_1(s, g)$递归地计算从状态$s$达成单个事实$g$的最小代价。

$h^{\max}$是可接纳的，但由于忽略了子目标之间的交互，通常过于乐观。

\subsection{$h^{\mathrm{add}}$启发式}

$h^{\mathrm{add}}$假设子目标之间完全独立，将所有子目标的代价相加：
$$h^{\mathrm{add}}(s) = \sum_{g \in G} h^{\mathrm{add}}_1(s, g)$$

$h^{\mathrm{add}}$通常比$h^{\max}$提供更好的搜索引导，但不是可接纳的（可能高估真实代价）。

\subsection{$h^{\mathrm{FF}}$启发式}

FF启发式（FastForward启发式）通过实际求解删除松弛问题来估计代价。其计算过程分两步：
\begin{enumerate}
  \item \textbf{前向传播：}在松弛规划图中逐层计算可达事实；
  \item \textbf{后向提取：}从目标出发反向提取松弛规划，其长度即为$h^{\mathrm{FF}}$值。
\end{enumerate}

$h^{\mathrm{FF}}$是FF规划器的核心，在实践中表现极为优秀。它既不是可接纳的也不是不可接纳的，但提供了极好的搜索引导。

\begin{appbox}[title={\textbf{应用：交通信号规划}}]
在城市交通信号控制中，可以将每个路口的信号状态建模为规划变量。$h^{\mathrm{FF}}$启发式可以高效地估计从当前信号配置到目标通行状态的调整步骤数，帮助快速求解大规模交通信号优化问题。
\end{appbox}

% =============================================
\section{规划图启发式 \ug}

\subsection{规划图（Planning Graph）}

规划图由Blum和Furst在1995年提出，是GraphPlan算法的核心数据结构。

\begin{definition}[规划图]
规划图是一个分层的有向图，交替包含：
\begin{itemize}
  \item \textbf{事实层（Fact Layer）：}包含在该层可能为真的命题；
  \item \textbf{动作层（Action Layer）：}包含前置条件在前一事实层中满足的动作。
\end{itemize}
每层还维护\textbf{互斥关系（Mutex）}，标记不能同时成立的事实或不能同时执行的动作。
\end{definition}

\subsection{GraphPlan算法}

GraphPlan交替进行两个阶段：
\begin{enumerate}
  \item \textbf{图扩展：}逐层构建规划图，直到所有目标事实出现在同一层且互不互斥；
  \item \textbf{解提取：}在规划图中后向搜索无互斥的动作集合。
\end{enumerate}

\subsection{$h^{\mathrm{level}}$启发式}

基于规划图的启发式函数，取目标事实全部首次出现且互不互斥的最早层号作为启发值。

% =============================================
\section{抽象启发式 \grad}

\subsection{模式数据库（Pattern Databases, PDB）}

\begin{definition}[模式数据库]
选取状态变量的一个子集（模式），将规划问题投影到该子集上形成更小的抽象问题。预先计算抽象问题中每个抽象状态到目标的精确距离，存储为查找表。
\end{definition}

PDB的关键挑战在于模式选择——如何选择"好"的变量子集使得启发值尽可能准确。

\subsection{合并与收缩（Merge-and-Shrink）}

合并与收缩是一种系统性构造抽象的方法：
\begin{enumerate}
  \item 从单个变量的转换系统出发；
  \item 迭代地合并（Merge）两个抽象系统并收缩（Shrink）以控制大小。
\end{enumerate}

\begin{keypoint}
合并与收缩产生的启发式是可接纳的和一致的，且在PDB启发式和$h^{\max}$之间建立了统一的理论框架。
\end{keypoint}

% =============================================
\section{地标启发式 \grad}

\subsection{地标（Landmarks）}

\begin{definition}[规划地标]
一个事实$l$是规划问题的地标，当且仅当在每个合法规划中，$l$至少在某个时刻为真。类似地，一个动作$a$是地标，当且仅当$a$出现在每个合法规划中。
\end{definition}

地标刻画了规划过程中的``必经之路''，提供了丰富的结构信息。

\subsection{LM-Cut启发式}

LM-Cut是一种基于地标的可接纳启发式，通过迭代地寻找代价最小割来计算启发值。它是目前最强的可接纳启发式之一。

\begin{algbox}[title={\textbf{LM-Cut核心思想}}]
\begin{enumerate}
  \item 计算$h^{\max}$代价，找到从初始状态到目标的最大代价路径上的一个``割''（cut）；
  \item 该割中动作的最小代价即为本轮启发值的贡献；
  \item 将割中动作的代价减去该最小值；
  \item 重复直到$h^{\max} = 0$，各轮贡献之和即为LM-Cut启发值。
\end{enumerate}
\end{algbox}

\subsection{地标计数启发式}

通过统计尚未实现的地标数量来估计剩余代价。虽然不是可接纳的，但计算简单且效果不错。

% =============================================
\section{其他启发式方法 \self}

\subsection{因果图启发式（Causal Graph Heuristic）}

利用变量之间的因果关系图，将多变量规划问题分解为一系列单变量问题逐个求解。

\subsection{势函数启发式（Potential Heuristics）}

将启发函数建模为状态变量取值的线性函数$h(s) = \sum_v w_{v,s(v)}$，通过线性规划优化权重$w$使得启发函数可接纳且尽可能精确。

\subsection{红黑规划启发式}

将变量分为``红色''（忽略删除效果）和``黑色''（完整处理），在两者之间取得平衡。

% =============================================
\section{本章小结与习题}

\subsection*{本章小结}
\begin{itemize}
  \item 删除松弛是规划领域最成功的启发式构造技术，衍生出$h^{\max}$、$h^{\mathrm{add}}$、$h^{\mathrm{FF}}$等重要启发式。
  \item 规划图提供了另一种启发式信息来源。
  \item 抽象启发式（PDB、合并与收缩）和地标启发式（LM-Cut）是最优规划的关键技术。
  \item 启发式设计的核心权衡：信息量 vs.\ 计算代价。
\end{itemize}

\subsection*{思考与练习}
\begin{enumerate}
  \item 对一个给定的规划问题，手动计算$h^{\max}$和$h^{\mathrm{add}}$值并比较。
  \item 构建一个简单问题的规划图并识别互斥关系。
  \item （研究生）证明LM-Cut启发式的可接纳性。
  \item （研究生）设计实验比较不同启发式在IPC基准问题上的求解能力。
\end{enumerate}

% 第3章：经典规划算法
\chapter{经典规划算法}

\begin{levelbox}
本章涵盖经典规划的核心算法。SAT/CSP规划（3.1--3.2节）为本科生内容，HTN与偏序规划（3.3--3.4节）为本科/研究生共同内容，时态与不确定性规划（3.5--3.6节）主要面向研究生。
\end{levelbox}

% =============================================
\section{基于SAT的规划 \ug}

\subsection{规划问题的SAT编码}

SAT规划的核心思想：将规划问题编码为命题逻辑公式，利用高效的SAT求解器来求解。

\begin{algbox}[title={\textbf{SAT规划基本框架}}]
\begin{enumerate}
  \item 设定一个时间步长上限$T$；
  \item 将规划问题编码为包含$T$个时间步的SAT公式$\phi_T$；
  \item 调用SAT求解器：若$\phi_T$可满足，从满足赋值中提取规划；若不可满足，增大$T$重试。
\end{enumerate}
\end{algbox}

编码需要处理以下约束：
\begin{itemize}
  \item 初始状态约束和目标约束；
  \item 动作前置条件和效果的因果约束；
  \item 帧公理（Frame Axioms）：未被动作改变的事实保持不变；
  \item 互斥约束：冲突动作不能同时执行。
\end{itemize}

\begin{appbox}[title={\textbf{应用：通信网络编码方案选择}}]
在通信系统中，需要根据信道条件、带宽需求和QoS约束选择合适的编码和传输方案。可以将可选的编码方式、调制方案和功率等级建模为SAT变量，将兼容性和性能要求建模为约束子句，利用SAT求解器自动找到满足所有约束的最优传输配置。
\end{appbox}

\subsection{SATplan与Madagascar}

\textbf{SATplan}是SAT规划的代表系统，多次在IPC竞赛中取得优异成绩。\textbf{Madagascar}进一步优化了编码效率和并行求解能力。

% =============================================
\section{基于CSP/SMT的规划 \ug}

\subsection{CSP规划}

约束满足问题（CSP）规划将规划建模为变量-值-约束的框架，适合处理复杂的资源和时间约束。

\subsection{SMT规划 \grad}

可满足性模理论（SMT）扩展了SAT，支持整数算术、实数算术等理论，能够处理数值变量和连续资源约束。

\begin{example}[SMT规划：资源调度]
考虑一个任务调度问题：$n$个任务需要分配到$m$台机器上，每个任务有处理时间、截止日期和资源需求。使用SMT编码可以自然地表达：
\begin{itemize}
  \item 时间变量：$t_i \in \mathbb{R}$（任务$i$的开始时间）
  \item 资源约束：任何时刻每台机器上的资源使用不超过容量
  \item 时序约束：前置任务必须先完成
\end{itemize}
\end{example}

% =============================================
\section{层次任务网络（HTN） \ug}

\subsection{HTN规划基本概念}

HTN（Hierarchical Task Network）规划采用自顶向下的任务分解策略，将复杂任务逐步分解为可直接执行的原子动作。

\begin{definition}[HTN规划]
一个HTN规划问题包括：
\begin{itemize}
  \item 原子任务（Primitive Tasks）：可直接执行的动作；
  \item 复合任务（Compound Tasks）：需要进一步分解的高层任务；
  \item 方法（Methods）：定义复合任务如何分解为子任务序列；
  \item 初始任务网络（Initial Task Network）：需要完成的顶层任务。
\end{itemize}
\end{definition}

\subsection{全序HTN与偏序HTN}

\begin{itemize}
  \item \textbf{全序HTN（Totally Ordered HTN）：}子任务之间有严格的线性顺序。SHOP/SHOP2是代表系统。
  \item \textbf{偏序HTN（Partially Ordered HTN）：}子任务之间只有部分顺序约束，允许并行执行。PANDA规划器支持此类问题。
\end{itemize}

\subsection{SHOP/SHOP2规划器}

SHOP（Simple Hierarchical Ordered Planner）是最著名的HTN规划器：
\begin{itemize}
  \item 使用全序分解，简化了搜索过程；
  \item 支持在前置条件中使用任意Lisp表达式；
  \item 在实际应用中表现优异，广泛用于军事规划、Web服务组合等领域。
\end{itemize}

\begin{appbox}[title={\textbf{应用：军事作战任务规划}}]
军事任务规划天然具有层次结构：战略目标分解为战役任务，战役任务再分解为具体的战术行动。HTN规划可以建模为：
\begin{itemize}
  \item 战略层：``夺取制空权''分解为``压制防空''$+$``空中巡逻''$+$``战场遮断''
  \item 战术层：``压制防空''分解为``电子干扰''$+$``反辐射打击''$+$``效果评估''
  \item 执行层：``电子干扰''分解为具体的频率选择和功率设置动作
\end{itemize}
\end{appbox}

% =============================================
\section{偏序规划（POP） \ug}

\subsection{偏序规划的基本思想}

偏序规划（Partial-Order Planning）在\textbf{规划空间}（而非状态空间）中搜索。它维护一个部分规划，通过逐步添加动作和排序约束来完善。

\begin{definition}[部分规划]
一个部分规划 $P = \langle A, \prec, \mathcal{B}, \mathcal{L} \rangle$ 包含：
\begin{itemize}
  \item $A$：动作集合；
  \item $\prec$：动作间的偏序关系；
  \item $\mathcal{B}$：变量绑定约束；
  \item $\mathcal{L}$：因果链接（Causal Links），记录哪个动作提供了哪个前置条件。
\end{itemize}
\end{definition}

\subsection{最小承诺策略}

偏序规划的核心原则是\textbf{最小承诺}（Least Commitment）：尽可能推迟决策，只在必要时添加排序约束，保留更多的灵活性。

\subsection{因果链接与威胁解决}

当一个动作可能破坏已建立的因果链接时，产生\textbf{威胁}（Threat）。解决威胁的方式有两种：
\begin{itemize}
  \item \textbf{提升（Promotion）：}将威胁动作排在因果链接之后；
  \item \textbf{降级（Demotion）：}将威胁动作排在因果链接之前。
\end{itemize}

% =============================================
\section{时态规划 \grad}

\subsection{持续性动作}

经典规划假设动作瞬时执行，而时态规划允许动作具有持续时间。PDDL 2.1引入了\textbf{持续性动作}（Durative Actions），包含：
\begin{itemize}
  \item 起始条件和效果（at start）
  \item 持续条件（over all）
  \item 终止条件和效果（at end）
  \item 持续时间约束
\end{itemize}

\subsection{时态规划器}

\begin{itemize}
  \item \textbf{TFD（Temporal Fast Downward）：}扩展Fast Downward支持时态规划，使用时态松弛的启发式。
  \item \textbf{OPTIC：}支持连续变化和时态约束的规划器。
  \item \textbf{Scotty：}使用凸优化处理连续时态规划问题。
\end{itemize}

\begin{appbox}[title={\textbf{应用：交通信号配时}}]
城市交通信号配时本质上是时态规划问题：每个信号相位有持续时间，不同方向的绿灯不能同时亮起，且需要满足最小绿灯时间、行人过街时间等约束。时态规划器可以自动生成满足安全和效率要求的信号配时方案。
\end{appbox}

% =============================================
\section{不确定性环境下的规划 \grad}

\subsection{一致性规划（Conformant Planning）}

在不确定初始状态下，一致性规划寻找一个在所有可能初始状态下都能达成目标的动作序列。本质上是在信念空间（Belief Space）中的搜索。

\subsection{条件规划（Contingent Planning）}

条件规划允许在执行过程中进行观测和分支：根据观测结果选择不同的后续动作。其输出是一棵\textbf{条件规划树}而非线性序列。

\subsection{MDP规划 \grad}

马尔可夫决策过程（MDP）以概率方式处理动作不确定性：

\begin{definition}[MDP]
一个MDP由$\langle S, A, T, R, \gamma \rangle$定义：
\begin{itemize}
  \item $S$：状态集合
  \item $A$：动作集合
  \item $T(s'|s,a)$：状态转移概率
  \item $R(s,a)$：回报函数
  \item $\gamma \in [0,1)$：折扣因子
\end{itemize}
\end{definition}

\textbf{值迭代}和\textbf{策略迭代}是求解MDP的经典算法。

\subsection{POMDP规划 \self}

部分可观测马尔可夫决策过程（POMDP）在MDP基础上增加了观测的不确定性。POMDP的精确求解是PSPACE-完全的，实践中通常使用近似方法（如PBVI、SARSOP）。

% =============================================
\section{本章小结与习题}

\subsection*{本章小结}
\begin{itemize}
  \item SAT/CSP/SMT将规划问题转化为约束求解问题，利用成熟的求解器技术。
  \item HTN通过任务分解捕获领域的层次结构，是实际应用中最成功的规划范式之一。
  \item 偏序规划的最小承诺策略保留了规划的灵活性。
  \item 时态规划和不确定性规划扩展了经典规划的表达能力。
\end{itemize}

\subsection*{思考与练习}
\begin{enumerate}
  \item 将一个3-任务调度问题编码为SAT公式并手动求解。
  \item 为``准备早餐''设计HTN方法层次。
  \item 比较前向状态空间搜索和偏序规划在积木世界问题上的搜索空间大小。
  \item （研究生）分析PDDL 2.1时态规划的表达能力与计算复杂性。
\end{enumerate}

% 第4章：现代经典规划器
\chapter{现代经典规划器 \ug}

\begin{levelbox}
本章介绍在国际规划竞赛（IPC）中表现优异的代表性规划器。掌握这些系统有助于理解规划算法的实际工程实现。本科生重点掌握FF和Fast Downward的基本原理，研究生应深入了解其优化技术。
\end{levelbox}

% =============================================
\section{FF规划器 \ug}

\subsection{FF的搜索框架}

FF（Fast-Forward）是Jörg Hoffmann于2001年开发的规划器，在IPC-2000竞赛中获得冠军。其核心技术包括：

\begin{enumerate}
  \item \textbf{$h^{\mathrm{FF}}$启发式：}通过求解删除松弛问题获得启发值（见第2章）；
  \item \textbf{有用动作（Helpful Actions）：}从松弛规划中提取``有用''的动作，优先搜索；
  \item \textbf{增强爬山搜索（Enforced Hill-Climbing, EHC）：}结合爬山和宽度优先搜索，在启发值改善时前进，陷入局部最优时用BFS逃脱。
\end{enumerate}

\begin{algbox}[title={\textbf{FF的搜索策略}}]
FF使用两级搜索框架：
\begin{enumerate}
  \item 首先尝试EHC：从当前状态出发，只保留使$h^{\mathrm{FF}}$严格降低的后继（限于有用动作），若无改善则用BFS寻找改善。
  \item 若EHC失败（无法找到降低$h$值的状态），切换到标准的GBFS搜索（使用完整动作集）。
\end{enumerate}
\end{algbox}

\subsection{FF的影响}

FF的设计思想深刻影响了后续规划器的开发：
\begin{itemize}
  \item 有用动作剪枝被广泛采用；
  \item $h^{\mathrm{FF}}$启发式成为基准参考；
  \item EHC搜索框架被多个后续系统继承。
\end{itemize}

% =============================================
\section{Fast Downward \ug}

\subsection{系统架构}

Fast Downward由Malte Helmert于2006年开发，是当前最具影响力的规划系统。其独特的三阶段架构：

\begin{enumerate}
  \item \textbf{翻译（Translation）：}将PDDL输入转换为紧凑的多值变量表示（SAS$^+$）；
  \item \textbf{知识编译（Knowledge Compilation）：}计算因果图、领域转换图等数据结构；
  \item \textbf{搜索（Search）：}使用各种启发式和搜索算法求解。
\end{enumerate}

\subsection{SAS$^+$表示}

与PDDL的命题表示不同，SAS$^+$使用有限域变量：每个变量$v$的取值域$D_v$可以有多个值。这种表示更紧凑，并能揭示问题的结构信息。

\subsection{Fast Downward的搜索组件}

Fast Downward支持丰富的搜索配置：
\begin{itemize}
  \item 搜索算法：A*、GBFS、加权A*、lazy GBFS等
  \item 启发函数：$h^{\mathrm{FF}}$、$h^{\mathrm{add}}$、因果图、LM-Cut、PDB、合并与收缩等
  \item 搜索增强：有用动作剪枝、preferred operators、多启发式open list
\end{itemize}

% =============================================
\section{LAMA规划器 \ug}

LAMA（Landmarks, Actions, and Multipath-search with $h^{\mathrm{add}}$）在IPC-2008和IPC-2011竞赛中获得冠军。

\subsection{LAMA的核心技术}
\begin{enumerate}
  \item \textbf{地标发现与利用：}自动发现规划地标，使用地标计数启发式；
  \item \textbf{多启发式搜索：}同时使用地标启发式和$h^{\mathrm{FF}}$，维护多个open list交替扩展；
  \item \textbf{迭代搜索：}找到第一个解后，用其代价作为上界继续搜索更优解。
\end{enumerate}

% =============================================
\section{Scorpion与其他现代规划器 \grad}

\subsection{Scorpion}

Scorpion基于饱和代价分区（Saturated Cost Partitioning），将多个抽象启发式的代价分配进行优化组合，在最优规划中取得了领先效果。

\subsection{Complementary系列}

基于互补启发式思想，使用在线选择机制动态决定每个搜索状态使用哪个启发函数。

\subsection{SymK}

SymK使用符号搜索（基于BDD——二叉决策图），可以紧凑表示和操作大规模状态集合。特别适合求解$K$个最优/准最优规划。

% =============================================
\section{国际规划竞赛（IPC）简介 \self}

\subsection{IPC的历史与赛道}

国际规划竞赛（International Planning Competition）始于1998年，每2--3年举办一次，推动了规划技术的持续进步。主要赛道包括：
\begin{itemize}
  \item \textbf{经典赛道（Classical Track）：}确定性、完全可观测的STRIPS/ADL规划
  \item \textbf{最优赛道（Optimal Track）：}要求找到代价最优的规划
  \item \textbf{满足赛道（Satisficing Track）：}只需找到合法规划，不要求最优
  \item \textbf{敏捷赛道（Agile Track）：}比拼求解速度
  \item \textbf{学习赛道（Learning Track）：}允许对训练问题进行离线学习
\end{itemize}

\subsection{IPC基准域}

IPC积累了丰富的基准规划域，包括：积木世界、物流、运输、卫星、电梯、供应链、网络安全等，是评估规划算法的标准测试集。

% =============================================
\section{本章小结与习题}

\subsection*{本章小结}
\begin{itemize}
  \item FF开创了基于松弛启发式和有用动作剪枝的高效满足规划范式。
  \item Fast Downward的模块化架构使其成为规划研究的标准平台。
  \item LAMA通过多启发式和地标技术在满足规划中取得突破。
  \item 现代最优规划器（Scorpion、SymK）利用抽象和符号方法求解更大规模问题。
\end{itemize}

\subsection*{思考与练习}
\begin{enumerate}
  \item 安装Fast Downward，用它求解一个简单的物流规划问题。
  \item 比较FF的EHC搜索和标准GBFS在同一问题上的节点扩展数。
  \item 解释LAMA为什么在满足规划竞赛中表现优于最优规划器。
  \item （研究生）阅读最近一届IPC的竞赛报告，分析获胜系统的关键技术。
\end{enumerate}


% =============================================
% 第二篇：元启发式与优化方法
% =============================================
\part{元启发式与优化方法}

% 第5章：元启发式与优化方法
\chapter{元启发式与优化方法}

\begin{levelbox}
元启发式方法在组合优化和任务调度中有广泛应用。遗传算法和模拟退火（5.1--5.2节）为本科生内容，群体智能算法（5.3节）为研究生内容，新型元启发式（5.4节）为自学拓展。
\end{levelbox}

% =============================================
\section{进化算法 \ug}

\subsection{遗传算法（GA）}

遗传算法模拟自然进化过程，通过选择、交叉和变异操作搜索解空间。

\begin{algbox}[title={\textbf{遗传算法基本框架}}]
\begin{enumerate}
  \item \textbf{编码：}将任务序列/资源分配方案编码为染色体（如排列编码、二进制编码）；
  \item \textbf{初始化：}随机生成初始种群；
  \item \textbf{适应度评估：}计算每个个体的目标函数值；
  \item \textbf{选择：}按适应度选择父代（轮盘赌、锦标赛选择等）；
  \item \textbf{交叉：}两个父代交换部分基因产生后代（顺序交叉OX、部分匹配交叉PMX等）；
  \item \textbf{变异：}以小概率改变基因值（交换变异、插入变异等）；
  \item \textbf{迭代：}重复直到满足终止条件。
\end{enumerate}
\end{algbox}

\textbf{在规划中的应用：}GA特别适合任务调度和资源分配问题。例如，Job-Shop调度问题中，将工序排列编码为染色体，使用最小完工时间（Makespan）作为适应度函数。

\begin{appbox}[title={\textbf{应用：车辆路径规划（VRP）}}]
在物流配送中，需要为一组车辆规划服务路线以最小化总行驶距离。GA可以有效求解大规模VRP问题：将每条路线编码为客户访问序列，使用特殊的交叉算子（如NWOX）保持路线的合法性。
\end{appbox}

\subsection{差分进化（DE）\grad}

差分进化使用种群中个体间的差异向量来驱动搜索：
$$\mathbf{v}_i = \mathbf{x}_{r_1} + F \cdot (\mathbf{x}_{r_2} - \mathbf{x}_{r_3})$$

DE在连续优化问题中表现优异，参数少且鲁棒性强。在连续域任务规划（如机器人轨迹优化）中有良好的应用。

% =============================================
\section{模拟退火与禁忌搜索 \ug}

\subsection{模拟退火（SA）}

\begin{algbox}[title={\textbf{模拟退火核心思想}}]
类比金属退火过程：
\begin{enumerate}
  \item 从初始解和初始高温$T_0$开始；
  \item 在当前解的邻域中随机选取新解$s'$；
  \item 若$f(s') < f(s)$（改善），接受$s'$；
  \item 若$f(s') \geq f(s)$（恶化），以概率$\exp(-\Delta f / T)$接受$s'$；
  \item 按冷却计划降低温度$T$；
  \item 重复直到温度足够低。
\end{enumerate}
\textbf{关键：}高温时容易接受较差解（探索），低温时趋向于只接受改善（开发）。
\end{algbox}

\subsection{禁忌搜索（TS）\grad}

禁忌搜索使用\textbf{禁忌表}（Tabu List）记录最近访问过的解（或移动），禁止短期内回访以跳出局部最优。

核心组件：
\begin{itemize}
  \item \textbf{禁忌表：}存储最近$k$步的移动，禁止这些移动的逆操作；
  \item \textbf{特赦准则（Aspiration Criteria）：}若某禁忌移动产生的解优于历史最优，则解除禁忌；
  \item \textbf{邻域结构：}定义从当前解到邻近解的转换操作。
\end{itemize}

% =============================================
\section{群体智能算法 \grad}

\subsection{粒子群优化（PSO）}

粒子群优化模拟鸟群觅食行为。每个粒子代表一个候选解，根据自身历史最优和群体历史最优更新位置：
$$\mathbf{v}_i^{t+1} = w\mathbf{v}_i^t + c_1 r_1 (\mathbf{p}_i - \mathbf{x}_i^t) + c_2 r_2 (\mathbf{g} - \mathbf{x}_i^t)$$
$$\mathbf{x}_i^{t+1} = \mathbf{x}_i^t + \mathbf{v}_i^{t+1}$$

\textbf{离散PSO：}标准PSO处理连续变量。用于组合优化时需要离散化处理，如将位置映射为排列（排列PSO）或使用集合操作定义速度更新。

\subsection{蚁群优化（ACO）}

蚁群优化模拟蚂蚁通过信息素寻找最短路径的行为：

\begin{enumerate}
  \item 每只蚂蚁根据信息素浓度和启发信息概率性地构建解；
  \item 构建完成后，根据解的质量更新信息素——好的解路径上信息素增加；
  \item 同时全局信息素以一定速率蒸发，避免过早收敛。
\end{enumerate}

ACO在TSP和路径规划中表现出色，其最优蚁群系统（MMAS）和蚁群系统（ACS）是两个重要变体。

\begin{appbox}[title={\textbf{应用：无人机航路规划}}]
ACO非常适合无人机航路规划：将航路点作为图的节点，信息素表示路径的优劣，启发信息包含距离、威胁和能耗等。蚂蚁从起点出发，概率性选择下一个航路点，最终找到安全高效的飞行路线。
\end{appbox}

% =============================================
\section{新型元启发式算法 \self}

近年来涌现了大量受自然现象启发的新型元启发式算法：

\subsection{灰狼优化（GWO）}

模拟灰狼群的社会等级和捕猎行为。狼群分为$\alpha$（最优）、$\beta$（次优）、$\delta$（第三）和$\omega$（其余），搜索过程模拟包围和攻击猎物。

\subsection{鲸鱼优化算法（WOA）}

模拟座头鲸的气泡网捕食策略，包含三种行为：包围猎物、螺旋游动和随机搜索。

\subsection{哈里斯鹰优化（HHO）}

模拟哈里斯鹰的协同捕猎策略，包含探索阶段和开发阶段（突袭行为）。

\begin{remark}
新型元启发式算法数量众多（100+种），但很多在本质上与已有算法相似。在选择使用时，应关注：(1) 是否有严格的理论分析；(2) 是否在标准基准上与经典方法进行了公平比较；(3) 是否解决了实际问题而非仅为新颖性。
\end{remark}

% =============================================
\section{元启发式在任务规划中的应用框架}

\subsection{编码策略}

将规划问题转化为优化问题的关键是设计合适的编码：
\begin{itemize}
  \item \textbf{排列编码：}适用于调度问题（任务执行顺序）
  \item \textbf{二进制编码：}适用于选择问题（是否执行某动作）
  \item \textbf{实数编码：}适用于连续参数优化（轨迹控制点）
  \item \textbf{混合编码：}同时包含离散和连续决策变量
\end{itemize}

\subsection{约束处理}

规划问题通常包含大量约束。常用的约束处理技术：
\begin{itemize}
  \item 罚函数法：将约束违反度加入目标函数
  \item 修复算子：对不可行解进行修复
  \item 解码器法：设计解码器保证总是生成可行解
\end{itemize}

% =============================================
\section{本章小结与习题}

\subsection*{本章小结}
\begin{itemize}
  \item 元启发式方法适合处理NP-hard的组合优化和调度问题，不依赖问题的特殊结构。
  \item GA和SA是最经典的元启发式，易于理解和实现。
  \item 群体智能算法（PSO、ACO）通过群体协作搜索，在路径规划等问题上有独特优势。
  \item 选择元启发式时应注重实效性而非新颖性。
\end{itemize}

\subsection*{思考与练习}
\begin{enumerate}
  \item 用遗传算法求解一个5任务3机器的Job-Shop调度问题。
  \item 比较模拟退火和禁忌搜索在TSP问题上的求解质量和收敛速度。
  \item （研究生）用ACO算法解决一个带时间窗的车辆路径问题（VRPTW）。
  \item （研究生）分析PSO和GWO在数学本质上的异同。
\end{enumerate}


% =============================================
% 第三篇：强化学习与规划
% =============================================
\part{强化学习与规划}

% 第6章：基于模型的强化学习与规划
\chapter{基于模型的强化学习与规划}

\begin{levelbox}
强化学习与规划的结合是当前AI的前沿方向。MDP/动态规划和蒙特卡洛树搜索（6.1--6.2节）为本科生核心内容，世界模型和AlphaGo/MuZero（6.3--6.4节）为研究生内容。
\end{levelbox}

% =============================================
\section{从MDP到强化学习规划 \ug}

\subsection{动态规划方法回顾}

在已知MDP模型的前提下，经典动态规划方法可以求解最优策略：

\textbf{值迭代（Value Iteration）：}
$$V_{k+1}(s) = \max_a \left[ R(s,a) + \gamma \sum_{s'} T(s'|s,a) V_k(s') \right]$$

\textbf{策略迭代（Policy Iteration）：}交替进行策略评估（计算$V^\pi$）和策略改进（贪心更新$\pi$）。

\subsection{Dyna架构 \ug}

Dyna（Sutton, 1991）是将学习与规划统一起来的开创性框架：

\begin{algbox}[title={\textbf{Dyna-Q核心思想}}]
\begin{enumerate}
  \item \textbf{与环境交互：}执行动作，获得真实经验$(s, a, r, s')$；
  \item \textbf{模型学习：}用真实经验更新环境模型$\hat{T}$和$\hat{R}$；
  \item \textbf{直接强化学习：}用真实经验更新$Q$值；
  \item \textbf{规划（模拟）：}从模型中采样$k$条模拟经验，用于额外的$Q$值更新。
\end{enumerate}
每次真实交互后进行$k$步规划，大幅提升样本效率。
\end{algbox}

% =============================================
\section{蒙特卡洛树搜索（MCTS） \ug}

\subsection{MCTS基本框架}

MCTS是一种基于采样的搜索方法，通过反复模拟来估计动作的价值。

\begin{algbox}[title={\textbf{MCTS四个阶段}}]
\begin{enumerate}
  \item \textbf{选择（Selection）：}从根节点出发，按照树策略（如UCT）选择子节点直到达到叶节点；
  \item \textbf{扩展（Expansion）：}在叶节点添加一个或多个新子节点；
  \item \textbf{模拟（Simulation）：}从新节点出发，使用默认策略（通常随机）进行快速模拟直到终态；
  \item \textbf{回溯（Backpropagation）：}将模拟结果沿路径回传更新所有祖先节点的统计信息。
\end{enumerate}
\end{algbox}

\subsection{UCT算法}

UCT（Upper Confidence Bounds for Trees）使用UCB1公式平衡探索与利用：
$$\mathrm{UCT}(s, a) = \bar{Q}(s, a) + c \sqrt{\frac{\ln N(s)}{N(s, a)}}$$
其中$\bar{Q}(s,a)$是动作$a$的平均回报，$N(s)$是状态$s$的访问次数，$N(s,a)$是动作$a$被选择的次数，$c$是探索常数。

\begin{appbox}[title={\textbf{应用：博弈规划}}]
MCTS在博弈中的应用最为经典。在围棋、国际象棋等游戏中，MCTS通过大量模拟评估每步棋的长期价值。MCTS的``随时''（Anytime）特性使其在限时决策中特别有用：可以随时返回当前最佳动作，给予更多时间则给出更好的结果。
\end{appbox}

% =============================================
\section{AlphaGo系列 \grad}

\subsection{AlphaGo（2016）}

AlphaGo首次击败围棋世界冠军，结合了四个关键组件：
\begin{enumerate}
  \item 基于人类棋谱训练的策略网络$p_\sigma$（监督学习）
  \item 通过自我对弈强化学习改进的策略网络$p_\rho$
  \item 预测游戏胜负的价值网络$v_\theta$
  \item 将神经网络与MCTS结合的搜索过程
\end{enumerate}

\subsection{AlphaGo Zero与AlphaZero（2017）}

\begin{keypoint}
AlphaGo Zero的突破：
\begin{itemize}
  \item \textbf{无需人类知识：}完全从随机开始自我对弈训练
  \item \textbf{统一的网络：}单一神经网络同时输出策略$\mathbf{p}$和价值$v$
  \item \textbf{简化的MCTS：}不使用快速走棋模拟，直接用价值网络评估叶节点
  \item \textbf{训练循环：}MCTS产生训练数据 $\rightarrow$ 训练网络 $\rightarrow$ 新网络指导MCTS
\end{itemize}
AlphaZero进一步推广到国际象棋和日本将棋，证明了方法的通用性。

国内方面，赵冬斌等在《控制理论与应用》上发表的``深度强化学习综述：兼论计算机围棋的发展''系统梳理了从AlphaGo到AlphaGo Zero的技术演进。腾讯开发的围棋AI``绝艺''（FineArt）也在世界计算机围棋大赛中多次获得冠军，是国内MCTS+深度学习在博弈规划中的重要实践。
\end{keypoint}

% =============================================
\section{世界模型与规划 \grad}

\subsection{MuZero（2020）}

MuZero是AlphaZero的重要推广，不需要已知的环境规则：

\begin{algbox}[title={\textbf{MuZero的三个学习函数}}]
\begin{enumerate}
  \item \textbf{表示函数$h$：}将观测映射到隐状态 $s^0 = h(o_1, \ldots, o_t)$
  \item \textbf{动力学函数$g$：}预测下一个隐状态和即时奖励 $(r^k, s^k) = g(s^{k-1}, a^k)$
  \item \textbf{预测函数$f$：}从隐状态预测策略和价值 $(\mathbf{p}^k, v^k) = f(s^k)$
\end{enumerate}
MuZero在隐状态空间中进行MCTS规划，无需学习完整的状态转移模型，只需学习对规划有用的信息。
\end{algbox}

\subsection{Dreamer系列 \grad}

Dreamer学习环境的世界模型，然后``在想象中''（in imagination）进行规划：

\begin{itemize}
  \item \textbf{DreamerV1（2020）：}使用循环状态空间模型（RSSM）学习世界模型，在潜在空间中用价值函数评估想象轨迹。
  \item \textbf{DreamerV2（2021）：}改用离散表示，在Atari游戏上首次超越无模型方法。
  \item \textbf{DreamerV3（2023）：}通过符号化分布（Symlog）预测和自适应调整机制，无需调参即可在多种任务上工作，包括Minecraft钻石挖掘。
\end{itemize}

\begin{keypoint}
世界模型规划的核心优势：
\begin{itemize}
  \item \textbf{样本效率：}通过模型生成大量模拟经验，减少与真实环境的交互次数。
  \item \textbf{可迁移性：}学到的世界模型可用于不同任务。
  \item \textbf{安全性：}可以在模型中``试错''而不影响真实环境。
\end{itemize}
核心挑战：模型误差的累积（Compounding Error）可能导致规划在长时间尺度上不可靠。
\end{keypoint}

% =============================================
\section{其他基于模型的RL规划方法 \self}

\subsection{MBPO（Model-Based Policy Optimization）}

使用短期模型rollout来扩充真实数据，通过控制rollout长度平衡模型精度和数据量。

\subsection{PlaNet（Planning Network）}

完全在隐状态空间中规划，使用CEM（交叉熵方法）进行在线规划。不训练策略网络，纯粹依赖规划。

\subsection{TD-MPC系列}

结合时序差分学习和模型预测控制（MPC），在连续控制任务中表现出色。TD-MPC2在多领域80+任务上展示了统一的单一世界模型。

% =============================================
\section{本章小结与习题}

\subsection*{本章小结}
\begin{itemize}
  \item Dyna架构统一了学习和规划，是基于模型RL的基础范式。
  \item MCTS通过采样和模拟评估动作价值，UCT公式平衡探索与利用。
  \item AlphaGo/AlphaZero/MuZero展示了MCTS+深度学习的强大能力。
  \item 世界模型（Dreamer系列）通过``在想象中规划''实现高样本效率。
\end{itemize}

\subsection*{思考与练习}
\begin{enumerate}
  \item 实现Dyna-Q算法并在简单迷宫问题上测试不同规划步数$k$的影响。
  \item 用MCTS求解一个简单的棋类游戏（如井字棋），比较不同模拟次数对决策质量的影响。
  \item （研究生）分析MuZero为什么不需要学习完整的状态重建。
  \item （研究生）讨论世界模型误差累积的原因及缓解策略。
\end{enumerate}

% 第7章：无模型强化学习与规划
\chapter{无模型强化学习与规划}

\begin{levelbox}
本章介绍不依赖环境模型的强化学习方法及其在规划中的应用。Q-learning和策略梯度（7.1--7.2节）为本科生基础，层次强化学习（7.3节）和序列决策Transformer（7.4节）为研究生内容。
\end{levelbox}

% =============================================
\section{值函数方法 \ug}

\subsection{Q-Learning}

Q-Learning是最基本的无模型强化学习算法，通过直接学习动作价值函数$Q(s,a)$来获得最优策略：
$$Q(s,a) \leftarrow Q(s,a) + \alpha \left[ r + \gamma \max_{a'} Q(s', a') - Q(s,a) \right]$$

\begin{remark}
关于强化学习方法的系统性综述，国内有多篇重要文献。高阳、陈世福和陆鑫在《自动化学报》上的``强化学习研究综述''是中文领域的经典入门文献；刘全等在《计算机学报》上的``深度强化学习综述''系统梳理了DQN以来的深度RL方法。此外，王雪松等在《自动化学报》上的``安全强化学习综述''讨论了RL在安全关键任务中的约束处理，对军事等高可靠性应用场景尤为重要。
\end{remark}

\subsection{DQN（Deep Q-Network） \ug}

DQN用深度神经网络近似$Q$函数，引入两个关键技术：
\begin{itemize}
  \item \textbf{经验回放（Experience Replay）：}将交互经验存储在缓冲区中，随机采样打破样本相关性。
  \item \textbf{目标网络（Target Network）：}使用定期更新的副本网络计算目标值，稳定训练过程。
\end{itemize}

\textbf{DQN的改进版本：}
\begin{itemize}
  \item \textbf{Double DQN：}解耦动作选择和价值评估，减少过高估计。
  \item \textbf{Dueling DQN：}分别估计状态价值$V(s)$和优势函数$A(s,a)$。
  \item \textbf{Prioritized Experience Replay：}优先采样TD误差大的经验。
  \item \textbf{Rainbow：}整合多种DQN改进技术。
\end{itemize}

% =============================================
\section{策略梯度方法 \ug}

\subsection{REINFORCE}

直接参数化策略$\pi_\theta(a|s)$，通过梯度上升最大化期望回报：
$$\nabla_\theta J(\theta) = \mathbb{E}_{\pi_\theta}\left[ \nabla_\theta \log \pi_\theta(a|s) \cdot G_t \right]$$

\subsection{Actor-Critic方法 \ug}

结合值函数和策略梯度：
\begin{itemize}
  \item \textbf{Actor（策略网络）：}决定采取什么动作
  \item \textbf{Critic（价值网络）：}评估动作的好坏，减小梯度估计的方差
\end{itemize}

\subsection{PPO（Proximal Policy Optimization） \grad}

PPO通过裁剪目标函数限制策略更新幅度：
$$L^{\mathrm{CLIP}}(\theta) = \mathbb{E}\left[\min\left(r_t(\theta)\hat{A}_t,\ \mathrm{clip}(r_t(\theta), 1-\epsilon, 1+\epsilon)\hat{A}_t\right)\right]$$
其中$r_t(\theta) = \pi_\theta(a_t|s_t)/\pi_{\theta_{\mathrm{old}}}(a_t|s_t)$。

PPO因其实现简单、效果稳定而成为最广泛使用的RL算法之一。

\subsection{SAC（Soft Actor-Critic） \grad}

SAC在最大化回报的同时最大化策略的熵，鼓励探索：
$$J(\pi) = \sum_t \mathbb{E}\left[ r(s_t, a_t) + \alpha \mathcal{H}(\pi(\cdot|s_t)) \right]$$

SAC在连续动作空间中的机器人控制任务中表现优异。

% =============================================
\section{层次强化学习 \grad}

\subsection{Options框架}

Options框架为RL引入了时间抽象，将原始动作扩展为``选项''（Option）：

\begin{definition}[Option]
一个选项$\omega = \langle I_\omega, \pi_\omega, \beta_\omega \rangle$包含：
\begin{itemize}
  \item $I_\omega \subseteq S$：启动集合（在哪些状态可以发起该选项）
  \item $\pi_\omega$：内部策略（选项执行期间的动作选择）
  \item $\beta_\omega: S \rightarrow [0,1]$：终止条件（在每个状态终止的概率）
\end{itemize}
\end{definition}

\subsection{HIRO（Hierarchical RL with Off-policy correction）}

HIRO使用目标条件层次结构：
\begin{itemize}
  \item 高层策略：设定子目标（goal）
  \item 低层策略：学习达成子目标的动作序列
  \item 通过离策略校正（off-policy correction）提升样本效率
\end{itemize}

\subsection{HAM和MaxQ}

\begin{itemize}
  \item \textbf{HAM（Hierarchy of Abstract Machines）：}使用有限状态机层次结构约束策略空间。
  \item \textbf{MaxQ：}将值函数分解为任务层次中各子任务的贡献之和，支持递归学习。
\end{itemize}

\begin{appbox}[title={\textbf{应用：军事任务分配中的层次决策}}]
军事作战中的任务分配天然具有层次结构。层次RL可以建模为：
\begin{itemize}
  \item 高层（指挥层）：分配战术目标给各作战单元
  \item 中层（协调层）：协调单元间的协同行动
  \item 低层（执行层）：控制具体的机动和打击动作
\end{itemize}
Options框架可以将``占领阵地''等复杂行动建模为时间扩展的选项。
\end{appbox}

% =============================================
\section{序列决策与Transformer \grad}

\subsection{Decision Transformer}

Decision Transformer将RL问题重新框定为序列建模问题：

\begin{keypoint}
核心思想：将RL轨迹视为序列 $(R_1, s_1, a_1, R_2, s_2, a_2, \ldots)$，其中$R_t$为``期望回报''（Return-to-Go）。使用GPT架构的Transformer建模该序列，推理时指定高$R_t$即可生成高质量的动作序列。

\textbf{优势：}
\begin{itemize}
  \item 无需值函数估计和时序差分学习
  \item 天然支持离线RL（从已有数据学习）
  \item 可以通过调节$R_t$控制行为质量
\end{itemize}
\end{keypoint}

\subsection{Trajectory Transformer \self}

将轨迹建模推向极致：使用Transformer同时建模状态转移和奖励，将规划实现为轨迹空间上的beam search。

\subsection{Gato \self}

DeepMind的Gato将多种任务（游戏、机器人控制、图像描述、对话）的数据统一为token序列，用单一Transformer处理，探索了通用智能体的可能性。

% =============================================
\section{本章小结与习题}

\subsection*{本章小结}
\begin{itemize}
  \item DQN将深度学习引入RL，经验回放和目标网络是关键技术。
  \item PPO和SAC是当前最实用的策略梯度算法。
  \item 层次RL通过时间抽象处理长时间尺度的规划问题。
  \item Decision Transformer将RL问题转化为序列建模，开辟了新的规划范式。
\end{itemize}

\subsection*{思考与练习}
\begin{enumerate}
  \item 实现Q-Learning并在简单的网格世界中训练智能体导航。
  \item 比较DQN和Double DQN在同一环境中的学习曲线。
  \item （研究生）分析Decision Transformer与传统RL方法在离线设定下的优劣。
  \item （研究生）设计一个层次RL框架来解决多步骤任务规划问题。
\end{enumerate}


% =============================================
% 第四篇：深度学习与规划
% =============================================
\part{深度学习与规划}

% 第8章：深度学习驱动的规划
\chapter{深度学习驱动的规划}

\begin{levelbox}
本章介绍深度学习在规划中的前沿应用。GNN规划（8.1节）和Transformer规划（8.2节）为研究生内容，扩散模型规划（8.3节）和神经符号规划（8.4节）为自学拓展。
\end{levelbox}

% =============================================
\section{图神经网络与规划 \grad}

\subsection{规划问题的图表示}

规划问题天然具有图结构：动作通过前置条件和效果与状态变量（命题）相连。将规划问题表示为图后，可以利用GNN来学习和泛化。

\subsection{ASNets（Action Schema Networks）}

ASNets是将GNN应用于规划的代表性工作：

\begin{algbox}[title={\textbf{ASNets核心思想}}]
\begin{enumerate}
  \item 将规划问题表示为二部图：动作节点和命题节点交替连接
  \item 使用消息传递网络在图上传播信息
  \item 网络参数在相同``类型''的动作/命题间共享（利用动作模式的对称性）
  \item 输出每个动作的执行概率（策略）或状态评估值（启发式）
\end{enumerate}
\textbf{泛化能力：}在小规模问题上训练，可以泛化到大规模同类问题。
\end{algbox}

\subsection{STRIPS-HGN（Hypergraph Networks）}

使用超图网络表示STRIPS问题，其中每个动作作为超边连接其前置条件和效果。超图结构更准确地捕获了动作与命题之间的多元关系。

% =============================================
\section{Transformer与规划 \grad}

\subsection{学习启发式函数}

使用Transformer从规划实例中学习启发式函数：
\begin{itemize}
  \item 将PDDL描述序列化为token序列
  \item Transformer编码器提取特征
  \item 输出层预测$h$值
  \item 可以在多个领域上联合训练，学习跨领域的规划知识
\end{itemize}

\subsection{Plansformer}

Plansformer直接使用Transformer端到端地生成规划：
\begin{enumerate}
  \item 输入：PDDL问题描述（序列化）
  \item 输出：动作序列
  \item 训练数据：由传统规划器生成的问题-规划对
\end{enumerate}

\subsection{基于注意力的搜索引导}

利用Transformer的注意力机制学习哪些状态特征在搜索中最重要，动态调整搜索策略。

% =============================================
\section{扩散模型与规划 \self}

\subsection{Diffuser}

Diffuser（Janner et al., 2022）将规划问题建模为轨迹去噪过程：

\begin{keypoint}
Diffuser核心思想：
\begin{enumerate}
  \item 将完整轨迹$\tau = (s_0, a_0, s_1, a_1, \ldots, s_T)$视为需要生成的对象
  \item 前向过程：逐步向轨迹添加噪声
  \item 反向过程：从纯噪声出发逐步去噪生成轨迹
  \item 条件生成：通过classifier-guided或classifier-free方式加入目标条件和约束
\end{enumerate}
\textbf{优势：}自然支持长期规划、多模态规划分布、灵活的约束引入。
\end{keypoint}

\subsection{Diffusion Policy}

Diffusion Policy将扩散模型用于机器人操作策略学习：
\begin{itemize}
  \item 以观测（图像、机器人状态）为条件，生成动作序列
  \item 相比其他生成模型，扩散模型更好地表达多模态动作分布
  \item 在多个机器人操作基准任务上取得了最佳效果
\end{itemize}

\subsection{扩散规划的优势与局限}

\textbf{优势：}
\begin{itemize}
  \item 天然支持多模态规划（同一情境下的多种可行方案）
  \item 通过引导（guidance）灵活加入各类约束和偏好
  \item 全局规划：同时生成整条轨迹，避免局部贪心
\end{itemize}

\textbf{局限：}
\begin{itemize}
  \item 推理速度较慢（需要多步去噪）
  \item 对训练数据的覆盖要求较高
  \item 理论分析尚不完善
\end{itemize}

% =============================================
\section{神经符号规划 \self}

\subsection{概述}

神经符号规划试图结合神经网络的感知学习能力和符号规划的推理保证：

\begin{itemize}
  \item \textbf{神经感知 + 符号规划：}用神经网络从原始输入（图像、传感器数据）提取符号表示，再用传统规划器求解。
  \item \textbf{可微规划：}将规划器嵌入可微分的计算图中，支持端到端训练。
  \item \textbf{神经PDDL：}自动从交互数据中学习PDDL领域模型。
\end{itemize}

\subsection{NLM（Neural Logic Machines）}

NLM使用张量化的逻辑运算实现可微分的逻辑推理：
\begin{itemize}
  \item 用矩阵表示谓词的真值
  \item 用可学习的卷积操作模拟逻辑规则
  \item 可以学习积木世界等领域的规划规则
\end{itemize}

\subsection{NSRT（Neuro-Symbolic Relational Transition models）}

NSRT学习可泛化的关系型状态转移模型：
\begin{itemize}
  \item 自动发现操作符（类似PDDL动作）
  \item 每个操作符关联一个采样器（用于生成连续参数）
  \item 支持Task and Motion Planning（TAMP）
\end{itemize}

% =============================================
\section{本章小结与习题}

\subsection*{本章小结}
\begin{itemize}
  \item GNN利用规划问题的图结构，实现跨问题规模的泛化。
  \item Transformer可以直接生成规划或学习启发式函数。
  \item 扩散模型将规划重新定义为轨迹生成问题，支持多模态和条件规划。
  \item 神经符号方法试图兼得神经网络的灵活性和符号推理的可靠性。
\end{itemize}

\subsection*{思考与练习}
\begin{enumerate}
  \item （研究生）分析GNN在规划中的泛化能力：为什么ASNets可以在更大规模问题上工作？
  \item （研究生）比较Diffuser和传统轨迹优化方法的优缺点。
  \item （研究生）讨论端到端神经规划和传统规划器在可靠性方面的差异。
  \item 思考：神经符号规划是否能真正结合两种范式的优点？存在哪些根本性挑战？
\end{enumerate}


% =============================================
% 第五篇：大语言模型与规划
% =============================================
\part{大语言模型与规划}

% 第9章：大语言模型与任务规划
\chapter{大语言模型与任务规划}

\begin{levelbox}
LLM规划是最活跃的研究前沿。LLM直接规划和混合方法（9.1--9.2节）为研究生核心内容，推理增强（9.3节）为研究生扩展内容，具身智能规划（9.4节）为自学拓展。
\end{levelbox}

% =============================================
\section{LLM直接规划 \grad}

\subsection{LLM作为规划器}

大语言模型（LLM）凭借海量知识和推理能力，可以直接生成任务规划：

\begin{algbox}[title={\textbf{LLM直接规划的基本方式}}]
\begin{itemize}
  \item \textbf{零样本规划：}直接向LLM描述任务和约束，让其生成动作序列
  \item \textbf{少样本规划：}提供几个规划示例作为上下文，引导LLM按格式输出
  \item \textbf{结构化输出：}要求LLM以PDDL、JSON等结构化格式输出规划
\end{itemize}
\end{algbox}

\subsection{LLM规划的优势与局限}

\textbf{优势：}
\begin{itemize}
  \item 丰富的世界知识：无需显式建模领域知识
  \item 自然语言接口：用户可以用自然语言描述任务
  \item 灵活性：可以处理开放域的新颖任务
\end{itemize}

\textbf{局限：}
\begin{itemize}
  \item \textbf{幻觉问题：}可能生成看似合理但实际不可行的规划
  \item \textbf{约束违反：}难以严格满足复杂的逻辑和资源约束
  \item \textbf{长序列规划：}随着步骤增加，错误累积导致规划质量下降
  \item \textbf{可验证性：}难以形式化验证规划的正确性
\end{itemize}

% =============================================
\section{LLM + 规划器混合方法 \grad}

\subsection{LLM+P框架}

LLM+P（Liu et al., 2023）将LLM的语言理解能力与经典规划器的求解能力结合：

\begin{algbox}[title={\textbf{LLM+P工作流程}}]
\begin{enumerate}
  \item LLM将自然语言问题描述翻译为PDDL格式
  \item 经典规划器（如Fast Downward）求解PDDL问题
  \item LLM将规划器输出翻译回自然语言解释
\end{enumerate}
\textbf{优势：}结合LLM的语言能力和规划器的求解保证。\\
\textbf{挑战：}PDDL翻译的准确性是瓶颈。
\end{algbox}

\subsection{LLM生成启发式}

一种新颖的结合方式：让LLM为特定领域生成启发式函数代码，然后在传统搜索框架中使用。

\begin{example}[LLM生成的启发式]
给定积木世界的问题描述，LLM可以生成如下Python启发式函数：
\begin{lstlisting}[language=Python]
def heuristic(state, goal):
    misplaced = 0
    for block, target_pos in goal.items():
        if state[block] != target_pos:
            misplaced += 1
    return misplaced
\end{lstlisting}
虽然这个启发式可能不是最优的，但LLM能够快速为新领域生成合理的启发式，无需人工设计。
\end{example}

\subsection{SayPlan \self}

SayPlan使用LLM进行大规模场景中的任务规划：
\begin{itemize}
  \item 使用3D场景图（3D Scene Graph）表示环境
  \item LLM通过迭代搜索场景图来定位相关物体和区域
  \item 结合场景约束验证来过滤不可行的规划步骤
\end{itemize}

% =============================================
\section{推理增强的LLM规划 \grad}

\subsection{思维链（Chain-of-Thought, CoT）}

\begin{keypoint}
CoT提示让LLM在给出最终答案前展示推理过程：
\begin{itemize}
  \item 分步推理减少跳跃性思维导致的错误
  \item 在规划中表现为：先分析当前状态、识别子目标、考虑约束，再给出动作
  \item 变体：Zero-shot CoT（``Let's think step by step''）、Auto-CoT等
\end{itemize}
\end{keypoint}

\subsection{思维树（Tree of Thoughts, ToT）}

ToT将线性推理扩展为树状搜索：
\begin{enumerate}
  \item 在每个决策点生成多个候选思路
  \item 评估每个思路的前景
  \item 使用BFS或DFS在思维树中搜索
  \item 支持回溯：放弃不佳的思路
\end{enumerate}

\subsection{思维图（Graph of Thoughts, GoT）\self}

GoT进一步将推理结构从树扩展为有向无环图（DAG），允许多个推理分支合并和交互。

\subsection{RAP（Reasoning via Planning）\self}

将LLM的推理过程建模为MDP，使用MCTS在推理空间中搜索，LLM既作为策略提供候选推理步骤，又作为世界模型评估状态价值。

% =============================================
\section{具身智能中的LLM规划 \self}

\subsection{SayCan}

\begin{algbox}[title={\textbf{SayCan：LLM + 可供性}}]
SayCan（Google, 2022）将LLM的语言理解与机器人的可供性（Affordance）结合：
\begin{enumerate}
  \item LLM生成候选动作的语言描述及其概率（``Say''）
  \item 预训练的技能策略评估每个动作在当前状态下的可行性（``Can''）
  \item 综合两者选择最佳动作：$\arg\max_a P_{\mathrm{LLM}}(a|q) \cdot P_{\mathrm{affordance}}(a|s)$
\end{enumerate}
\end{algbox}

\subsection{Code as Policies}

直接让LLM生成机器人控制代码（Python），利用预定义的API与机器人交互。LLM的代码生成能力使其可以处理空间推理、循环和条件等复杂逻辑。

\subsection{VoxPoser \self}

使用LLM和视觉-语言模型在3D体素空间中生成约束和可供性地图，指导机器人的运动规划。零样本地将语言指令转化为机器人动作。

\subsection{RT-2与PaLM-E \self}

\begin{itemize}
  \item \textbf{RT-2（Robotic Transformer 2）：}将视觉-语言模型直接用作机器人策略，输出离散化的动作token。
  \item \textbf{PaLM-E：}具身多模态语言模型，将传感器数据（图像、点云）作为token嵌入语言模型，统一处理语言理解和机器人规划。
\end{itemize}

\begin{appbox}[title={\textbf{应用：智能仓储中的任务规划}}]
在智能仓储场景中，LLM可以理解自然语言的订单描述，并为机器人集群生成取货-打包-配送的任务序列。SayCan式的方法可以确保规划的每一步都是机器人物理上可执行的。
\end{appbox}

% =============================================
\section{本章小结与习题}

\subsection*{本章小结}
\begin{itemize}
  \item LLM直接规划灵活但缺乏保证；LLM+规划器混合方法兼顾两者优势。
  \item CoT/ToT/GoT通过结构化推理提升LLM的规划能力。
  \item 具身智能中的LLM规划（SayCan、Code as Policies等）是机器人应用的热点方向。
  \item LLM规划的核心挑战是可靠性和可验证性。
\end{itemize}

\subsection*{思考与练习}
\begin{enumerate}
  \item 使用LLM API为一个简单的家庭任务（如准备咖啡）生成规划，分析其合理性。
  \item （研究生）比较LLM+P和LLM直接规划在同一问题集上的成功率。
  \item （研究生）设计一个实验评估CoT和ToT在任务规划中的效果差异。
  \item 讨论：LLM规划何时可以替代传统规划器？何时不能？
\end{enumerate}

% 第10章：语言智能体
\chapter{语言智能体 \grad}

\begin{levelbox}
语言智能体是LLM规划的前沿应用方向。ReAct和Reflexion（10.1--10.2节）为研究生核心内容，多智能体LLM和自主探索智能体（10.3--10.4节）为自学拓展。
\end{levelbox}

% =============================================
\section{ReAct：推理与行动的协同 \grad}

\subsection{ReAct框架}

ReAct（Yao et al., 2023）将推理（Reasoning）和行动（Acting）交织在一起：

\begin{algbox}[title={\textbf{ReAct循环}}]
\begin{enumerate}
  \item \textbf{思考（Thought）：}LLM分析当前状况，推理下一步该做什么
  \item \textbf{行动（Action）：}LLM生成一个具体动作（如搜索、查询、操作）
  \item \textbf{观测（Observation）：}从环境获取动作的执行结果
  \item 重复以上循环直到任务完成
\end{enumerate}
与CoT的区别：ReAct能够与外部环境交互获取真实信息，而非仅凭LLM的内部知识推理。
\end{algbox}

\subsection{ReAct在任务规划中的应用}

ReAct特别适合需要信息收集和动态决策的规划任务：
\begin{itemize}
  \item 在线购物：搜索商品、比较价格、下单
  \item 信息检索：多步搜索和推理回答复杂问题
  \item 软件操作：在Web界面上完成多步操作
\end{itemize}

% =============================================
\section{Reflexion：自我反思的规划 \grad}

\subsection{Reflexion框架}

Reflexion（Shinn et al., 2023）让智能体从失败中学习：

\begin{algbox}[title={\textbf{Reflexion工作流程}}]
\begin{enumerate}
  \item \textbf{执行（Act）：}智能体尝试完成任务
  \item \textbf{评估（Evaluate）：}判断任务是否成功
  \item \textbf{反思（Reflect）：}若失败，LLM生成对失败原因的文字分析
  \item \textbf{重试（Retry）：}将反思结果加入上下文，重新尝试
\end{enumerate}
\textbf{关键：}反思是以自然语言形式存储在记忆中的经验教训，类似人类的``经验总结''。
\end{algbox}

\subsection{LATS（Language Agent Tree Search）\self}

LATS将MCTS引入语言智能体：
\begin{itemize}
  \item 将ReAct轨迹视为搜索树的路径
  \item 使用LLM自评估作为节点价值估计
  \item 环境反馈作为评价信号
  \item 支持回溯和多路径探索
\end{itemize}

% =============================================
\section{多智能体LLM系统 \self}

\subsection{AutoGen}

AutoGen（Microsoft）是一个多智能体对话框架：
\begin{itemize}
  \item 支持多个LLM智能体之间的对话和协作
  \item 智能体可以扮演不同角色（程序员、评审者、用户代理等）
  \item 通过对话协议编排复杂的工作流
\end{itemize}

\subsection{MetaGPT}

MetaGPT为LLM智能体引入标准化操作流程（SOP）：
\begin{itemize}
  \item 模拟软件公司的角色分工：产品经理、架构师、程序员、测试员
  \item 使用结构化文档（需求文档、设计文档、代码审查）进行信息传递
  \item 通过角色间的协作自动完成软件开发
\end{itemize}

\subsection{ChatDev}

ChatDev让多个LLM角色通过对话协作开发软件：
\begin{itemize}
  \item 模拟``聊天驱动''的软件开发流程
  \item 各角色依次参与设计、编码、测试、文档等环节
  \item 展示了LLM多智能体在创造性任务中的协同能力
\end{itemize}

% =============================================
\section{自主探索与终身学习智能体 \self}

\subsection{Voyager}

Voyager（Wang et al., 2023）是在Minecraft中的终身学习智能体：

\begin{keypoint}
Voyager的三个核心组件：
\begin{enumerate}
  \item \textbf{自动课程（Automatic Curriculum）：}LLM根据当前能力和探索进度自动设定下一个学习目标
  \item \textbf{技能库（Skill Library）：}将已学会的行为以代码形式存储，可复用和组合
  \item \textbf{迭代提示（Iterative Prompting）：}根据环境反馈和执行错误迭代改进代码
\end{enumerate}
Voyager展示了LLM驱动的开放世界持续探索和能力增长。
\end{keypoint}

\subsection{DEPS（Describe, Explain, Plan, Select）\self}

DEPS为LLM智能体提供了系统的错误恢复机制：先描述当前状况，解释失败原因，重新规划，选择最佳方案。

\subsection{Generative Agents \self}

Generative Agents（Park et al., 2023）在虚拟小镇中模拟了25个智能体的社会行为：
\begin{itemize}
  \item 每个智能体有记忆流、反思和规划模块
  \item 智能体可以形成社交关系、传播信息、协调活动
  \item 展示了LLM在模拟复杂社会行为方面的潜力
\end{itemize}

% =============================================
\section{本章小结与习题}

\subsection*{本章小结}
\begin{itemize}
  \item ReAct通过推理-行动交替实现了LLM与环境的有效交互。
  \item Reflexion通过自我反思实现了从失败中学习。
  \item 多智能体LLM系统（AutoGen、MetaGPT）探索了角色协作的规划模式。
  \item Voyager等系统展示了LLM驱动的开放世界终身学习智能体的可能性。
\end{itemize}

\subsection*{思考与练习}
\begin{enumerate}
  \item （研究生）实现一个简化的ReAct框架，让LLM通过工具使用完成一个多步骤任务。
  \item （研究生）比较ReAct和Reflexion在同一任务上的成功率和效率。
  \item 讨论：多智能体LLM系统中的``涌现行为''是真正的智能还是表面模式？
  \item 分析Voyager技能库的设计：为什么代码形式比自然语言形式更适合存储技能？
\end{enumerate}


% =============================================
% 第六篇：多智能体规划
% =============================================
\part{多智能体规划}

% 第11章：多智能体规划
\chapter{多智能体规划}

\begin{levelbox}
多智能体规划涉及多个智能体的协调与协作。MAPF基础和博弈论方法（11.1--11.2节）为本科/研究生内容，多智能体RL（11.3节）为研究生内容，分布式规划（11.4节）为自学拓展。
\end{levelbox}

% =============================================
\section{多智能体路径规划（MAPF） \ug}

\subsection{问题定义}

\begin{definition}[MAPF问题]
给定一个图$G=(V,E)$和$k$个智能体，每个智能体$i$有起始位置$s_i$和目标位置$g_i$。要求为所有智能体规划无冲突的路径（同一时刻不能有两个智能体占据同一节点或沿同一边反向移动）。
\end{definition}

\textbf{优化目标：}
\begin{itemize}
  \item \textbf{Makespan：}最小化所有智能体完成任务的最大时间
  \item \textbf{Sum of Costs：}最小化所有智能体路径代价之和
\end{itemize}

\subsection{CBS（Conflict-Based Search） \ug}

\begin{algbox}[title={\textbf{CBS算法}}]
CBS是一种双层搜索算法：
\begin{enumerate}
  \item \textbf{低层搜索：}为每个智能体独立规划路径（遵守给定的约束）
  \item \textbf{高层搜索：}在约束树（Constraint Tree）中搜索
    \begin{itemize}
      \item 根节点：无约束，每个智能体走自己的最短路径
      \item 检测冲突：两个智能体在同一时刻占据同一位置
      \item 分支：对冲突生成两个子节点，分别约束其中一个智能体避开冲突位置/时刻
      \item 重复直到找到无冲突的方案
    \end{itemize}
\end{enumerate}
\textbf{最优性：}CBS在Sum of Costs目标下返回最优解。
\end{algbox}

\subsection{ICTS（Increasing Cost Tree Search） \grad}

ICTS也是双层搜索：
\begin{itemize}
  \item 高层搜索：递增代价树，枚举各智能体路径代价的可能组合
  \item 低层搜索：验证给定代价组合下是否存在无冲突方案
\end{itemize}

\subsection{LaCAM与LaCAM* \grad}

LaCAM是一种高效的亚最优MAPF算法：
\begin{itemize}
  \item 在联合配置空间中进行搜索，但使用惰性约束生成避免指数爆炸
  \item LaCAM*是其改进版本，在保持效率的同时提供更好的解质量
  \item 可以在秒级求解数千个智能体的MAPF问题
\end{itemize}

\begin{appbox}[title={\textbf{应用：自动驾驶交叉路口协调}}]
在城市交叉路口，多辆自动驾驶车辆需要在避免碰撞的前提下高效通过。MAPF算法可以建模为：每辆车是一个智能体，路口的离散网格/车道为图的节点，需要在时空中规划无冲突的通行方案。
\end{appbox}

% =============================================
\section{博弈论方法 \grad}

\subsection{博弈论与多智能体规划}

当多个智能体有不同甚至对立的目标时，需要博弈论框架来分析和求解。

\begin{definition}[正则形式博弈]
一个$n$人博弈由 $\langle N, (S_i)_{i \in N}, (u_i)_{i \in N} \rangle$ 定义：
\begin{itemize}
  \item $N = \{1, \ldots, n\}$：参与者集合
  \item $S_i$：参与者$i$的策略集合
  \item $u_i: S_1 \times \cdots \times S_n \rightarrow \mathbb{R}$：参与者$i$的效用函数
\end{itemize}
\end{definition}

\subsection{纳什均衡与规划}

纳什均衡是博弈的基本解概念：没有任何参与者可以通过单方面改变策略来提高自己的效用。在多智能体规划中，纳什均衡可以作为多个自利智能体交互的预测结果。

\subsection{Stackelberg博弈}

Stackelberg博弈描述了领导者-跟随者的层次结构：
\begin{itemize}
  \item 领导者先选择策略
  \item 跟随者观察到领导者的策略后做出最优响应
  \item 在安全博弈、资源部署等领域有重要应用
\end{itemize}

% =============================================
\section{多智能体强化学习（MARL） \grad}

\subsection{MARL基本框架}

多智能体环境通常建模为\textbf{随机博弈}（Stochastic Game）或\textbf{Dec-POMDP}（分布式POMDP）。

\begin{remark}
MARL在国内是活跃的研究方向。梁星星、冯旸赫等在《自动化学报》上的``多Agent深度强化学习综述''系统梳理了MARL的算法体系和应用场景。陈杰、方浩和辛斌的专著《多智能体系统的协同群集运动控制》和向峥嵘等的《多智能体协同控制及应用》则从控制论视角讨论了多智能体协同的理论与方法。
\end{remark}

\subsection{CTDE范式（集中训练分布式执行）}

\begin{keypoint}
CTDE（Centralized Training with Decentralized Execution）是MARL的主流范式：
\begin{itemize}
  \item \textbf{训练时：}所有智能体的信息集中处理，学习全局最优的协作策略
  \item \textbf{执行时：}每个智能体仅基于自己的局部观测做出决策
\end{itemize}
\end{keypoint}

\subsection{QMIX}

QMIX将联合动作价值函数分解为各智能体的个体价值函数的单调混合：
$$Q_{\mathrm{tot}}(\boldsymbol{\tau}, \mathbf{a}) = f_{\mathrm{mix}}(Q_1(\tau_1, a_1), \ldots, Q_n(\tau_n, a_n); s)$$
其中混合网络$f_{\mathrm{mix}}$的权重均为非负，保证单调性。

\subsection{MAPPO}

MAPPO将PPO扩展到多智能体设定，每个智能体有独立的Actor和共享的Critic：
\begin{itemize}
  \item Actor：基于局部观测输出动作
  \item Critic：基于全局状态评估价值
  \item 实现简单，在StarCraft II等基准上效果优异
\end{itemize}

\begin{appbox}[title={\textbf{应用：无人机蜂群协同}}]
无人机蜂群执行搜索救援、区域监视等任务时，需要多架无人机协同规划。MARL可以学习：
\begin{itemize}
  \item 区域分配：各无人机负责不同区域
  \item 编队控制：保持通信连通性的同时覆盖最大区域
  \item 任务协调：动态分配新发现的目标
\end{itemize}
QMIX或MAPPO可以在模拟环境中训练协同策略，然后部署到真实无人机上。

国内在无人机蜂群协同规划方面有大量研究。王建峰、贾高伟等在《系统工程与电子技术》上的综述系统分析了多无人机协同任务规划方法。张宏宏、李文华等在《航空学报》上的``有人/无人机协同作战：概念、技术与挑战''则从军事应用视角讨论了协同规划中的关键技术问题。
\end{appbox}

% =============================================
\section{分布式规划 \self}

\subsection{MA-STRIPS}

MA-STRIPS将经典STRIPS规划扩展到多智能体设定：
\begin{itemize}
  \item 每个智能体有自己的动作集合和局部视角
  \item 通过通信协调生成联合规划
  \item 隐私保护：智能体不需要完全共享自己的内部信息
\end{itemize}

\subsection{Plan Merging}

规划合并方法先让每个智能体独立生成自己的局部规划，然后通过协商合并为一致的联合规划：
\begin{enumerate}
  \item 各智能体独立规划
  \item 检测规划间的冲突
  \item 通过协商解决冲突（重排序、等待、替代动作）
\end{enumerate}

\subsection{联盟形成}

联盟形成（Coalition Formation）研究智能体如何动态组成合作小组以完成任务：
\begin{itemize}
  \item 任务分配与联盟组建的联合优化
  \item 基于能力匹配和代价分摊的联盟稳定性分析
  \item 在军事多平台协同中有重要应用
\end{itemize}

% =============================================
\section{本章小结与习题}

\subsection*{本章小结}
\begin{itemize}
  \item MAPF是多智能体规划的基础问题，CBS是最优求解的代表算法。
  \item 博弈论为利益冲突的多智能体交互提供分析框架。
  \item MARL（QMIX、MAPPO）通过集中训练分布式执行实现多智能体协作规划。
  \item 分布式规划保护智能体隐私的同时实现协调。
\end{itemize}

\subsection*{思考与练习}
\begin{enumerate}
  \item 在一个$8 \times 8$的网格上，用CBS为4个智能体规划无冲突路径。
  \item 分析CBS在最坏情况下的计算复杂度。
  \item （研究生）比较QMIX和MAPPO在一个协同导航任务中的学习效率。
  \item （研究生）设计一个场景说明纳什均衡可能不是社会最优的。
  \item 讨论：在军事协同规划中，集中式和分布式规划各有什么优缺点？
\end{enumerate}


% =============================================
% 附录
% =============================================
\appendix
% 附录A：算法分类速查表与学习路线图
\chapter{算法分类速查表与学习路线图}

本附录汇总全书涉及的主要算法，提供分类速查和学习建议。

% =============================================
\section*{A.1 算法总表}

\begin{longtable}{p{3.5cm}p{4cm}p{1.2cm}p{4.5cm}}
\toprule
\textbf{算法类别} & \textbf{代表算法} & \textbf{层次} & \textbf{核心特点} \\
\midrule
\endfirsthead
\toprule
\textbf{算法类别} & \textbf{代表算法} & \textbf{层次} & \textbf{核心特点} \\
\midrule
\endhead
\bottomrule
\endfoot

\multicolumn{4}{l}{\textbf{第一篇：经典规划方法}} \\
\midrule
无信息搜索 & BFS, DFS, IDS & \textcolor{ugcolor}{本} & 不使用启发信息 \\
启发式搜索 & A*, IDA*, GBFS & \textcolor{ugcolor}{本} & 启发函数引导搜索 \\
加权A* & WA* & \textcolor{gradcolor}{研} & 牺牲最优性换速度 \\
松弛启发式 & $h^{\max}$, $h^{\mathrm{add}}$, $h^{\mathrm{FF}}$ & \textcolor{ugcolor}{本} & 删除松弛估计 \\
规划图 & GraphPlan & \textcolor{ugcolor}{本} & 互斥分析+解提取 \\
抽象启发式 & PDB, Merge\&Shrink & \textcolor{gradcolor}{研} & 问题投影和抽象 \\
地标启发式 & LM-Cut, 地标计数 & \textcolor{gradcolor}{研} & 必经之路分析 \\
SAT规划 & SATplan, Madagascar & \textcolor{ugcolor}{本} & 转化为SAT求解 \\
CSP/SMT规划 & CSP编码, SMT规划 & \textcolor{ugcolor}{本}/\textcolor{gradcolor}{研} & 约束满足求解 \\
HTN规划 & SHOP/SHOP2, PANDA & \textcolor{ugcolor}{本} & 层次任务分解 \\
偏序规划 & POP & \textcolor{ugcolor}{本} & 最小承诺策略 \\
时态规划 & TFD, OPTIC & \textcolor{gradcolor}{研} & 处理持续性动作 \\
MDP/POMDP & 值迭代, PBVI & \textcolor{gradcolor}{研} & 概率不确定性 \\
FF规划器 & FF (EHC+$h^{\mathrm{FF}}$) & \textcolor{ugcolor}{本} & 松弛启发+爬山 \\
Fast Downward & FD (SAS$^+$) & \textcolor{ugcolor}{本} & 模块化框架 \\
LAMA & 地标+多启发式 & \textcolor{ugcolor}{本} & 满足规划冠军 \\
Scorpion/SymK & 饱和代价/符号搜索 & \textcolor{gradcolor}{研} & 最优规划前沿 \\
\midrule

\multicolumn{4}{l}{\textbf{第二篇：元启发式与优化方法}} \\
\midrule
遗传算法 & GA & \textcolor{ugcolor}{本} & 选择交叉变异 \\
模拟退火 & SA & \textcolor{ugcolor}{本} & 概率接受+冷却 \\
差分进化 & DE & \textcolor{gradcolor}{研} & 差异向量驱动 \\
禁忌搜索 & TS & \textcolor{gradcolor}{研} & 禁忌表+特赦 \\
粒子群优化 & PSO & \textcolor{gradcolor}{研} & 群体信息共享 \\
蚁群优化 & ACO, MMAS & \textcolor{gradcolor}{研} & 信息素引导 \\
新型元启发式 & GWO, WOA, HHO & \textcolor{selfcolor}{拓} & 自然现象启发 \\
\midrule

\multicolumn{4}{l}{\textbf{第三篇：强化学习与规划}} \\
\midrule
Dyna架构 & Dyna-Q & \textcolor{ugcolor}{本} & 学习+规划统一 \\
蒙特卡洛树搜索 & MCTS, UCT & \textcolor{ugcolor}{本} & 采样+模拟评估 \\
AlphaGo/Zero & MCTS+深度网络 & \textcolor{gradcolor}{研} & 自我对弈+搜索 \\
MuZero & 学习的世界模型 & \textcolor{gradcolor}{研} & 隐空间规划 \\
Dreamer系列 & DreamerV1--V3 & \textcolor{gradcolor}{研} & 想象中的规划 \\
DQN系列 & DQN, Rainbow & \textcolor{ugcolor}{本} & 深度Q网络 \\
策略梯度 & REINFORCE, A2C & \textcolor{ugcolor}{本} & 直接优化策略 \\
PPO & PPO & \textcolor{gradcolor}{研} & 裁剪策略更新 \\
SAC & SAC & \textcolor{gradcolor}{研} & 最大熵RL \\
层次RL & Options, HIRO & \textcolor{gradcolor}{研} & 时间抽象 \\
Decision Transformer & DT & \textcolor{gradcolor}{研} & 序列建模=RL \\
\midrule

\multicolumn{4}{l}{\textbf{第四篇：深度学习与规划}} \\
\midrule
图神经网络规划 & ASNets & \textcolor{gradcolor}{研} & 图结构泛化 \\
Transformer规划 & Plansformer & \textcolor{gradcolor}{研} & 端到端生成 \\
扩散模型规划 & Diffuser, Diff.\ Policy & \textcolor{selfcolor}{拓} & 轨迹去噪生成 \\
神经符号规划 & NLM, NSRT & \textcolor{selfcolor}{拓} & 感知+推理融合 \\
\midrule

\multicolumn{4}{l}{\textbf{第五篇：大语言模型与规划}} \\
\midrule
LLM直接规划 & Zero/Few-shot规划 & \textcolor{gradcolor}{研} & 语言知识驱动 \\
LLM+规划器 & LLM+P & \textcolor{gradcolor}{研} & 混合求解 \\
思维链/树/图 & CoT, ToT, GoT & \textcolor{gradcolor}{研} & 结构化推理 \\
具身LLM规划 & SayCan, Code as Policies & \textcolor{selfcolor}{拓} & 机器人+LLM \\
语言智能体 & ReAct, Reflexion & \textcolor{gradcolor}{研} & 推理+行动+反思 \\
多智能体LLM & AutoGen, MetaGPT & \textcolor{selfcolor}{拓} & 角色协作 \\
自主探索 & Voyager & \textcolor{selfcolor}{拓} & 终身学习 \\
\midrule

\multicolumn{4}{l}{\textbf{第六篇：多智能体规划}} \\
\midrule
MAPF & CBS, ICTS & \textcolor{ugcolor}{本}/\textcolor{gradcolor}{研} & 无冲突路径 \\
LaCAM & LaCAM/LaCAM* & \textcolor{gradcolor}{研} & 大规模MAPF \\
博弈论方法 & Nash, Stackelberg & \textcolor{gradcolor}{研} & 利益冲突建模 \\
MARL & QMIX, MAPPO & \textcolor{gradcolor}{研} & 集中训练分布执行 \\
分布式规划 & MA-STRIPS & \textcolor{selfcolor}{拓} & 隐私保护协调 \\
\end{longtable}

% =============================================
\section*{A.2 推荐学习路线}

\subsection*{本科生核心路线（建议32--48学时）}

\begin{enumerate}
  \item \textbf{基础（8--12学时）：}第1章（搜索基础）$\rightarrow$ 第2章1--2节（松弛启发式、规划图）
  \item \textbf{经典算法（8--12学时）：}第3章1--4节（SAT、HTN、POP）$\rightarrow$ 第4章1--3节（FF、FD、LAMA）
  \item \textbf{优化方法（4--6学时）：}第5章1--2节（GA、SA）
  \item \textbf{强化学习（8--12学时）：}第6章1--2节（Dyna、MCTS）$\rightarrow$ 第7章1--2节（DQN、策略梯度）
  \item \textbf{多智能体（4--6学时）：}第11章1节（MAPF、CBS）
\end{enumerate}

\subsection*{研究生扩展路线（在本科基础上增加16--32学时）}

\begin{enumerate}
  \item \textbf{高级启发式（4--6学时）：}第2章3--4节（PDB、LM-Cut）
  \item \textbf{时态与不确定性（4--6学时）：}第3章5--6节
  \item \textbf{高级RL（4--8学时）：}第6章3--4节（AlphaZero、Dreamer）$\rightarrow$ 第7章3--4节（层次RL、DT）
  \item \textbf{深度学习规划（4--6学时）：}第8章1--2节（GNN、Transformer）
  \item \textbf{LLM规划（4--8学时）：}第9章1--3节 $\rightarrow$ 第10章1--2节
  \item \textbf{多智能体进阶（4--6学时）：}第11章2--3节（博弈论、MARL）
\end{enumerate}

\subsection*{自学拓展专题}

以下内容适合感兴趣的读者根据研究方向选择阅读：
\begin{itemize}
  \item 扩散模型与规划（第8章3节）
  \item 神经符号规划（第8章4节）
  \item 具身智能LLM规划（第9章4节）
  \item 语言智能体前沿（第10章3--4节）
  \item 分布式规划（第11章4节）
  \item 新型元启发式综述（第5章4节）
\end{itemize}

% =============================================
\section*{A.3 算法演化脉络图}

任务规划算法的发展可以概括为以下几条主线：

\begin{enumerate}
  \item \textbf{符号推理主线：}STRIPS $\rightarrow$ GraphPlan $\rightarrow$ SAT规划 $\rightarrow$ 启发式搜索（FF/FD） $\rightarrow$ 地标/抽象启发式
  \item \textbf{层次分解主线：}HTN $\rightarrow$ SHOP/SHOP2 $\rightarrow$ PANDA $\rightarrow$ 时态HTN
  \item \textbf{优化求解主线：}SA/GA $\rightarrow$ PSO/ACO $\rightarrow$ 混合元启发式
  \item \textbf{学习规划主线：}Dyna $\rightarrow$ DQN $\rightarrow$ AlphaGo/MuZero $\rightarrow$ Dreamer $\rightarrow$ Decision Transformer
  \item \textbf{深度生成主线：}VAE/GAN $\rightarrow$ Transformer $\rightarrow$ 扩散模型 $\rightarrow$ 轨迹生成规划
  \item \textbf{语言智能主线：}GPT $\rightarrow$ CoT/ToT $\rightarrow$ ReAct/Reflexion $\rightarrow$ Voyager $\rightarrow$ 多智能体LLM
\end{enumerate}

这些主线并非孤立发展，而是在不断交叉融合。例如，MuZero结合了学习和搜索，LLM+P结合了语言模型和符号规划，ASNets结合了深度学习和经典规划。未来的发展趋势是各主线的进一步融合，形成更强大的统一规划框架。

% =============================================
\section*{A.4 国内推荐参考文献}

\subsection*{教材与专著}
\begin{enumerate}
  \item 蔡自兴, 刘丽珏, 蔡竞峰.《人工智能及其应用》（第7版）. 清华大学出版社, 2024. ——国内使用最广泛的AI教材之一，含智能规划章节。
  \item 吴飞, 潘云鹤.《人工智能引论》. 高等教育出版社, 2024. ——``101计划''核心教材，体现国内AI教育前沿。
  \item Ghallab等著, 殷建平等译.《自动规划：理论和实践》. 清华大学出版社, 2008. ——国际经典教材中文版，系统讲述自动规划理论。
  \item 段海滨.《蚁群算法原理及其应用》. 科学出版社, 2005. ——蚁群算法的理论与工程应用。
  \item 焦李成等.《免疫优化计算、学习与识别》. 科学出版社, 2006. ——免疫计算与优化理论。
  \item 焦李成等.《深度学习、优化与识别》. 清华大学出版社, 2017. ——深度学习与优化方法。
  \item 陈杰, 辛斌.《智能优化的探索--开发权衡理论与方法》. 科学出版社, 2013. ——元启发式核心理论。
  \item 陈杰, 方浩, 辛斌.《多智能体系统的协同群集运动控制》. 科学出版社, 2017. ——多智能体协同控制。
  \item 向峥嵘等.《多智能体协同控制及应用》. 化学工业出版社, 2023. ——多智能体系统理论与应用。
  \item 余伶俐, 周开军, 陈白帆.《智能驾驶技术：路径规划与导航控制》. 机械工业出版社, 2020. ——自动驾驶中的规划技术。
  \item 赵鑫等.《大语言模型》. 电子工业出版社, 2024. ——大语言模型技术全面参考。
\end{enumerate}

\subsection*{重要综述论文}
\begin{enumerate}
  \item 高阳, 陈世福, 陆鑫. 强化学习研究综述. 自动化学报, 2004. ——中文RL经典入门综述。
  \item 刘全等. 深度强化学习综述. 计算机学报, 2018. ——深度RL方法体系梳理。
  \item 赵冬斌等. 深度强化学习综述：兼论计算机围棋的发展. 控制理论与应用, 2016. ——AlphaGo技术解析。
  \item 梁星星, 冯旸赫等. 多Agent深度强化学习综述. 自动化学报, 2020. ——MARL方法全面综述。
  \item 杜永浩, 邢立宁, 蔡昭权. 无人飞行器集群智能调度技术综述. 自动化学报, 2020. ——无人机调度方法体系。
  \item 贾高伟, 王建峰. 无人机集群任务规划方法研究综述. 系统工程与电子技术, 2021. ——无人机集群规划综述。
  \item 邢立宁, 陈英武. 任务规划系统研究及发展. 火力与指挥控制, 2006. ——任务规划系统发展脉络。
  \item 秦龙等. 基于大语言模型的复杂任务自主规划处理框架. 自动化学报, 2024. ——LLM任务规划前沿。
  \item 张宏宏等. 有人/无人机协同作战：概念、技术与挑战. 航空学报, 2024. ——有人/无人协同综述。
  \item 王建峰等. 多无人机协同任务规划方法研究综述. 系统工程与电子技术, 2024. ——多无人机协同最新进展。
  \item 王雪松等. 安全强化学习综述. 自动化学报, 2023. ——安全RL约束方法。
\end{enumerate}


% ============ 参考文献 ============
\begin{thebibliography}{99}
\addcontentsline{toc}{chapter}{参考文献}

\bibitem{ghallab2004} Ghallab M, Nau D, Traverso P. Automated Planning: Theory and Practice[M]. Morgan Kaufmann, 2004.
\bibitem{ghallab2016} Ghallab M, Nau D, Traverso P. Automated Planning and Acting[M]. Cambridge University Press, 2016.
\bibitem{russell2020} Russell S, Norvig P. Artificial Intelligence: A Modern Approach (4th ed)[M]. Pearson, 2020.
\bibitem{lavalle2006} LaValle S M. Planning Algorithms[M]. Cambridge University Press, 2006.
\bibitem{sutton2018} Sutton R S, Barto A G. Reinforcement Learning: An Introduction (2nd ed)[M]. MIT Press, 2018.
\bibitem{shoham2009} Shoham Y, Leyton-Brown K. Multiagent Systems: Algorithmic, Game-Theoretic, and Logical Foundations[M]. Cambridge University Press, 2009.

\bibitem{hoffmann2001} Hoffmann J, Nebel B. The FF Planning System: Fast Plan Generation Through Heuristic Search[J]. JAIR, 2001, 14:253--302.
\bibitem{helmert2006} Helmert M. The Fast Downward Planning System[J]. JAIR, 2006, 26:191--246.
\bibitem{blum1997} Blum A L, Furst M L. Fast Planning Through Planning Graph Analysis[J]. Artificial Intelligence, 1997, 90(1-2):281--300.
\bibitem{kautz1992} Kautz H, Selman B. Planning as Satisfiability[C]. ECAI, 1992.
\bibitem{erol1994} Erol K, Hendler J, Nau D S. HTN Planning: Complexity and Expressivity[C]. AAAI, 1994.
\bibitem{fox2003} Fox M, Long D. PDDL2.1: An Extension to PDDL for Expressing Temporal Planning Domains[J]. JAIR, 2003, 20:61--124.

\bibitem{silver2016} Silver D, et al. Mastering the Game of Go with Deep Neural Networks and Tree Search[J]. Nature, 2016, 529(7587):484--489.
\bibitem{silver2017} Silver D, et al. Mastering the Game of Go without Human Knowledge[J]. Nature, 2017, 550(7676):354--359.
\bibitem{schrittwieser2020} Schrittwieser J, et al. Mastering Atari, Go, Chess and Shogi by Planning with a Learned Model[J]. Nature, 2020, 584(7839):186--190.
\bibitem{hafner2023} Hafner D, et al. Mastering Diverse Domains through World Models[J]. arXiv:2301.04104, 2023.
\bibitem{chen2021} Chen L, et al. Decision Transformer: Reinforcement Learning via Sequence Modeling[C]. NeurIPS, 2021.

\bibitem{janner2022} Janner M, Du Y, Tenenbaum J B, Levine S. Planning with Diffusion for Flexible Behavior Synthesis[C]. ICML, 2022.
\bibitem{chi2023} Chi C, et al. Diffusion Policy: Visuomotor Policy Learning via Action Diffusion[C]. RSS, 2023.

\bibitem{ahn2022} Ahn M, et al. Do As I Can, Not As I Say: Grounding Language in Robotic Affordances[C]. CoRL, 2022.
\bibitem{yao2023} Yao S, et al. ReAct: Synergizing Reasoning and Acting in Language Models[C]. ICLR, 2023.
\bibitem{wei2022} Wei J, et al. Chain-of-Thought Prompting Elicits Reasoning in Large Language Models[C]. NeurIPS, 2022.
\bibitem{yao2023tot} Yao S, et al. Tree of Thoughts: Deliberate Problem Solving with Large Language Models[C]. NeurIPS, 2023.
\bibitem{liu2023llmp} Liu B, et al. LLM+P: Empowering Large Language Models with Optimal Planning Proficiency[J]. arXiv:2304.11477, 2023.
\bibitem{wang2023voyager} Wang G, et al. Voyager: An Open-Ended Embodied Agent with Large Language Models[J]. arXiv:2305.16291, 2023.

\bibitem{sharon2015} Sharon G, et al. Conflict-Based Search for Optimal Multi-Agent Pathfinding[J]. Artificial Intelligence, 2015, 219:40--66.
\bibitem{stern2019} Stern R, et al. Multi-Agent Pathfinding: Definitions, Variants, and Benchmarks[C]. SoCS, 2019.

% ============ 中文教材与专著 ============

\bibitem{cai2024} 蔡自兴, 刘丽珏, 蔡竞峰. 人工智能及其应用（第7版）[M]. 北京: 清华大学出版社, 2024.
\bibitem{wufei2024} 吴飞, 潘云鹤. 人工智能引论[M]. 北京: 高等教育出版社, 2024.
\bibitem{yin2008} Ghallab M, Nau D, Traverso P. 自动规划：理论和实践[M]. 殷建平, 等译. 北京: 清华大学出版社, 2008.
\bibitem{chenjie2017} 陈杰, 方浩, 辛斌. 多智能体系统的协同群集运动控制[M]. 北京: 科学出版社, 2017.
\bibitem{duanhaibin2005} 段海滨. 蚁群算法原理及其应用[M]. 北京: 科学出版社, 2005.
\bibitem{jiao2017} 焦李成, 赵进, 杨淑媛, 等. 深度学习、优化与识别[M]. 北京: 清华大学出版社, 2017.
\bibitem{jiao2006} 焦李成, 杜海峰, 刘芳, 等. 免疫优化计算、学习与识别[M]. 北京: 科学出版社, 2006.
\bibitem{chenjie2013} 陈杰, 辛斌. 智能优化的探索--开发权衡理论与方法[M]. 北京: 科学出版社, 2013.
\bibitem{yu2020} 余伶俐, 周开军, 陈白帆. 智能驾驶技术：路径规划与导航控制[M]. 北京: 机械工业出版社, 2020.
\bibitem{xiang2023} 向峥嵘, 等. 多智能体协同控制及应用[M]. 北京: 化学工业出版社, 2023.
\bibitem{zhaoxin2024} 赵鑫, 等. 大语言模型[M]. 北京: 电子工业出版社, 2024.

% ============ 中文综述论文 ============

\bibitem{gaoyang2004} 高阳, 陈世福, 陆鑫. 强化学习研究综述[J]. 自动化学报, 2004, 30(1):86--100.
\bibitem{liuquan2018} 刘全, 翟建伟, 章宗长, 等. 深度强化学习综述[J]. 计算机学报, 2018, 41(1):1--27.
\bibitem{zhaodongbin2016} 赵冬斌, 邵坤, 朱圆恒, 等. 深度强化学习综述: 兼论计算机围棋的发展[J]. 控制理论与应用, 2016, 33(6):701--717.
\bibitem{liang2020} 梁星星, 冯旸赫, 马跃, 等. 多Agent深度强化学习综述[J]. 自动化学报, 2020, 46(12):2537--2557.
\bibitem{du2020} 杜永浩, 邢立宁, 蔡昭权. 无人飞行器集群智能调度技术综述[J]. 自动化学报, 2020, 46(2):222--241.
\bibitem{jia2021} 贾高伟, 王建峰. 无人机集群任务规划方法研究综述[J]. 系统工程与电子技术, 2021, 43(1):99--111.
\bibitem{xing2006} 邢立宁, 陈英武. 任务规划系统研究及发展[J]. 火力与指挥控制, 2006, 31(4):1--4.
\bibitem{qin2024} 秦龙, 武万森, 等. 基于大语言模型的复杂任务自主规划处理框架[J]. 自动化学报, 2024, 50(4):862--872.
\bibitem{zhang2024} 张宏宏, 李文华, 郑家毅, 等. 有人/无人机协同作战: 概念、技术与挑战[J]. 航空学报, 2024, 45(15):29653.
\bibitem{wang2024} 王建峰, 贾高伟, 郭正, 等. 多无人机协同任务规划方法研究综述[J]. 系统工程与电子技术, 2024, 46(10):3437--3450.
\bibitem{wangxs2023} 王雪松, 王荣荣, 程玉虎. 安全强化学习综述[J]. 自动化学报, 2023, 49(3):489--511.

\end{thebibliography}

\end{document}
